\begin{filecontents*}[overwrite]{references.bib}
@article{batzner2022,
  author = {Batzner, Simon and Musaelian, Albert and Sun, Lixin and Geiger, Mario and Mailoa, Jonathan P. and Kornbluth, Mordechai and Molinari, Nicola and Smidt, Tess E. and Kozinsky, Boris},
  title = {E(3)-equivariant graph neural networks for data-efficient and accurate interatomic potentials},
  journal = {Nat. Commun.},
  year = {2022},
  volume = {13},
  pages = {2453},
  doi = {10.1038/s41467-022-29939-5},
}

@article{chen2022m3gnet,
  author = {Chen, Chi and Ong, Shyue Ping},
  title = {A universal graph deep learning interatomic potential for the periodic table},
  journal = {Nat. Comput. Sci.},
  year = {2022},
  volume = {2},
  pages = {718--728},
  doi = {10.1038/s43588-022-00349-3},
}

@article{musaelian2023allegro,
  author = {Musaelian, Albert and Batzner, Simon and Johansson, Anders and Sun, Lixin and Owen, Cameron J. and Kornbluth, Mordechai and Kozinsky, Boris},
  title = {Learning local equivariant representations for large-scale atomistic dynamics},
  journal = {Nat. Commun.},
  year = {2023},
  volume = {14},
  pages = {579},
  doi = {10.1038/s41467-023-36329-y},
}

@article{deng2023chgnet,
  author = {Deng, Bowen and Zhong, Peichen and Jun, KyuJung and Riebesell, Janosh and Han, Kevin and Bartel, Christopher J. and Ceder, Gerbrand},
  title = {CHGNet as a pretrained universal neural network potential for charge-informed atomistic modelling},
  journal = {Nat. Mach. Intell.},
  year = {2023},
  volume = {5},
  pages = {1031--1041},
  doi = {10.1038/s42256-023-00716-3},
}

@article{merchant2023gnome,
  author = {Merchant, Amil and Batzner, Simon and Schoenholz, Samuel S. and Aykol, Muratahan and Cheon, Gowoon and Cubuk, Ekin Dogus},
  title = {Scaling deep learning for materials discovery},
  journal = {Nature},
  year = {2023},
  volume = {624},
  pages = {80--85},
  doi = {10.1038/s41586-023-06735-9},
}

@article{batatia2025macemp0,
  author = {Batatia, Ilyes and Benner, Philipp and Chiang, Yuan and Elena, Alin M. and Kovács, Dávid P. and Riebesell, Janosh and Advincula, Xavier R. and Asta, Mark and Avaylon, Matthew and Baldwin, William J. and Berger, Fabian and Bernstein, Noam and Bhowmik, Arghya and Bigi, Filippo and Blau, Samuel M. and Cărare, Vlad and Ceriotti, Michele and Chong, Sanggyu and Darby, James P. and De, Sandip and Della Pia, Flaviano and Deringer, Volker L. and Elijošius, Rokas and El-Machachi, Zakariya and Fako, Edvin and Falcioni, Fabio and Ferrari, Andrea C. and Gardner, John L. A. and Gawkowski, Mikołaj J. and Genreith-Schriever, Annalena and George, Janine and Goodall, Rhys E. A. and Grandel, Jonas and Grey, Clare P. and Grigorev, Petr and Han, Shuang and Handley, Will and Heenen, Hendrik H. and Hermansson, Kersti and Ho, Cheuk Hin and Hofmann, Stephan and Holm, Christian and Jaafar, Jad and Jakob, Konstantin S. and Jung, Hyunwook and Kapil, Venkat and Kaplan, Aaron D. and Karimitari, Nima and Kermode, James R. and Kourtis, Panagiotis and Kroupa, Namu and Kullgren, Jolla and Kuner, Matthew C. and Kuryla, Domantas and Liepuoniute, Guoda and Lin, Chen and Margraf, Johannes T. and Magdău, Ioan-Bogdan and Michaelides, Angelos and Moore, J. Harry and Naik, Aakash A. and Niblett, Samuel P. and Norwood, Sam Walton and O’Neill, Niamh and Ortner, Christoph and Persson, Kristin A. and Reuter, Karsten and Rosen, Andrew S. and Rosset, Louise A. M. and Schaaf, Lars L. and Schran, Christoph and Shi, Benjamin X. and Sivonxay, Eric and Stenczel, Tamás K. and Sutton, Christopher and Svahn, Viktor and Swinburne, Thomas D. and Tilly, Jules and van der Oord, Cas and Vargas, Santiago and Varga-Umbrich, Eszter and Vegge, Tejs and Vondrák, Martin and Wang, Yangshuai and Witt, William C. and Wolf, Thomas and Zills, Fabian and Csányi, Gábor},
  title = {A foundation model for atomistic materials chemistry},
  journal = {J. Chem. Phys.},
  year = {2025},
  volume = {163},
  pages = {184110},
  doi = {10.1063/5.0297006},
}

@article{riebesell2025matbench,
  author = {Riebesell, Janosh and Goodall, Rhys E. A. and Benner, Philipp and Chiang, Yuan and Deng, Bowen and Ceder, Gerbrand and Asta, Mark and Lee, Alpha A. and Jain, Anubhav and Persson, Kristin A.},
  title = {A framework to evaluate machine learning crystal stability predictions},
  journal = {Nat. Mach. Intell.},
  year = {2025},
  volume = {7},
  pages = {836--847},
  doi = {10.1038/s42256-025-01055-1},
}

@article{deng2025softening,
  author = {Deng, Bowen and Choi, Yunyeong and Zhong, Peichen and Riebesell, Janosh and Anand, Shashwat and Li, Zhuohan and Jun, KyuJung and Persson, Kristin A. and Ceder, Gerbrand},
  title = {Systematic softening in universal machine learning interatomic potentials},
  journal = {npj Comput. Mater.},
  year = {2025},
  volume = {11},
  pages = {9},
  doi = {10.1038/s41524-024-01500-6},
}

@article{focassio2025surfaces,
  author = {Focassio, Bruno and M. Freitas, Luis Paulo and Schleder, Gabriel R.},
  title = {Performance Assessment of Universal Machine Learning Interatomic Potentials: Challenges and Directions for Materials’ Surfaces},
  journal = {ACS Appl. Mater. Interfaces},
  year = {2025},
  volume = {17},
  pages = {13111--13121},
  doi = {10.1021/acsami.4c03815},
}

@article{radova2025frozen,
  author = {Radova, Mariia and Stark, Wojciech G. and Allen, Connor S. and Maurer, Reinhard J. and Bartók, Albert P.},
  title = {Fine-tuning foundation models of materials interatomic potentials with frozen transfer learning},
  journal = {npj Comput. Mater.},
  year = {2025},
  volume = {11},
  pages = {237},
  doi = {10.1038/s41524-025-01727-x},
}

@article{kim2025reewc,
  author = {Kim, Jisu and Lee, Jiho and Oh, Sangmin and Park, Yutack and Hwang, Seungwoo and Han, Seungwu and Kang, Sungwoo and Kang, Youngho},
  title = {An efficient forgetting-aware fine-tuning framework for pretrained universal machine-learning interatomic potentials},
  journal = {npj Comput. Mater.},
  year = {2026},
  volume = {12},
  pages = {26},
  doi = {10.1038/s41524-025-01895-w},
}

@inproceedings{dong2024survey,
  author = {Dong, Qingxiu and Li, Lei and Dai, Damai and Zheng, Ce and Ma, Jingyuan and Li, Rui and Xia, Heming and Xu, Jingjing and Wu, Zhiyong and Chang, Baobao and Sun, Xu and Li, Lei and Sui, Zhifang},
  title = {A Survey on In-context Learning},
  booktitle = {Proc. 2024 Conf. Empirical Methods in Natural Language Processing},
  year = {2024},
  pages = {1107--1128},
  doi = {10.18653/v1/2024.emnlp-main.64},
}

@article{hollmann2025tabpfn,
  author = {Hollmann, Noah and Müller, Samuel and Purucker, Lennart and Krishnakumar, Arjun and Körfer, Max and Hoo, Shi Bin and Schirrmeister, Robin Tibor and Hutter, Frank},
  title = {Accurate predictions on small data with a tabular foundation model},
  journal = {Nature},
  year = {2025},
  volume = {637},
  pages = {319--326},
  doi = {10.1038/s41586-024-08328-6},
}

@article{jablonka2024predictive,
  author = {Jablonka, Kevin Maik and Schwaller, Philippe and Ortega-Guerrero, Andres and Smit, Berend},
  title = {Leveraging large language models for predictive chemistry},
  journal = {Nat. Mach. Intell.},
  year = {2024},
  volume = {6},
  pages = {161--169},
  doi = {10.1038/s42256-023-00788-1},
}

@article{jablonka2023examples,
  author = {Jablonka, Kevin Maik and Ai, Qianxiang and Al-Feghali, Alexander and Badhwar, Shruti and Bocarsly, Joshua D. and Bran, Andres M. and Bringuier, Stefan and Brinson, L. Catherine and Choudhary, Kamal and Circi, Defne and Cox, Sam and de Jong, Wibe A. and Evans, Matthew L. and Gastellu, Nicolas and Genzling, Jerome and Gil, María Victoria and Gupta, Ankur K. and Hong, Zhi and Imran, Alishba and Kruschwitz, Sabine and Labarre, Anne and Lála, Jakub and Liu, Tao and Ma, Steven and Majumdar, Sauradeep and Merz, Garrett W. and Moitessier, Nicolas and Moubarak, Elias and Mouriño, Beatriz and Pelkie, Brenden and Pieler, Michael and Ramos, Mayk Caldas and Ranković, Bojana and Rodriques, Samuel G. and Sanders, Jacob N. and Schwaller, Philippe and Schwarting, Marcus and Shi, Jiale and Smit, Berend and Smith, Ben E. and Van Herck, Joren and Völker, Christoph and Ward, Logan and Warren, Sean and Weiser, Benjamin and Zhang, Sylvester and Zhang, Xiaoqi and Zia, Ghezal Ahmad and Scourtas, Aristana and Schmidt, K. J. and Foster, Ian and White, Andrew D. and Blaiszik, Ben},
  title = {14 examples of how LLMs can transform materials science and chemistry: a reflection on a large language model hackathon},
  journal = {Digit. Discovery},
  year = {2023},
  volume = {2},
  pages = {1233--1250},
  doi = {10.1039/d3dd00113j},
}

@article{zheng2023mof,
  author = {Zheng, Zhiling and Zhang, Oufan and Borgs, Christian and Chayes, Jennifer T. and Yaghi, Omar M.},
  title = {ChatGPT Chemistry Assistant for Text Mining and the Prediction of MOF Synthesis},
  journal = {J. Am. Chem. Soc.},
  year = {2023},
  volume = {145},
  pages = {18048--18062},
  doi = {10.1021/jacs.3c05819},
}

@article{ramos2026boicl,
  author = {Caldas Ramos, Mayk and Michtavy, Shane S. and White, Andrew D. and Porosoff, Marc D.},
  title = {Bayesian Optimization of Catalysis with In-Context Learning},
  journal = {ACS Cent. Sci.},
  year = {2026},
  volume = {12},
  pages = {599--615},
  doi = {10.1021/acscentsci.5c02418},
}

@article{ekuma2025dynamicicl,
  author = {Ekuma, Chinedu E.},
  title = {Dynamic in-context learning with conversational models for data extraction and materials property prediction},
  journal = {APL Mach. Learn.},
  year = {2025},
  volume = {3},
  pages = {016119},
  doi = {10.1063/5.0254406},
}

@inproceedings{wang2025crystalicl,
  author = {Wang, Ruobing and Tan, Qiaoyu and Wang, Yili and Wang, Ying and Wang, Xin},
  title = {CrystalICL: Enabling In-Context Learning for Crystal Generation},
  booktitle = {Proc. 2025 Conf. Empirical Methods in Natural Language Processing},
  year = {2025},
  pages = {18440--18455},
  doi = {10.18653/v1/2025.emnlp-main.929},
}

@article{romera2024funsearch,
  author = {Romera-Paredes, Bernardino and Barekatain, Mohammadamin and Novikov, Alexander and Balog, Matej and Kumar, M. Pawan and Dupont, Emilien and Ruiz, Francisco J. R. and Ellenberg, Jordan S. and Wang, Pengming and Fawzi, Omar and Kohli, Pushmeet and Fawzi, Alhussein},
  title = {Mathematical discoveries from program search with large language models},
  journal = {Nature},
  year = {2024},
  volume = {625},
  pages = {468--475},
  doi = {10.1038/s41586-023-06924-6},
}

@inproceedings{shojaee2025llmsr,
  author = {Shojaee, Parshin and Meidani, Kazem and Gupta, Shashank and Farimani, Amir Barati and Reddy, Chandan K.},
  title = {{LLM-SR}: Scientific Equation Discovery via Programming with Large Language Models},
  booktitle = {The Thirteenth International Conference on Learning Representations (ICLR)},
  year = {2025},
  note = {arXiv:2404.18400},
}

@inproceedings{merler2024icsr,
  author = {Merler, Matteo and Haitsiukevich, Katsiaryna and Dainese, Nicola and Marttinen, Pekka},
  title = {In-Context Symbolic Regression: Leveraging Large Language Models for Function Discovery},
  booktitle = {Proc. 62nd Annual Meeting of the Association for Computational Linguistics (Volume 4: Student Research Workshop)},
  year = {2024},
  pages = {589--606},
  doi = {10.18653/v1/2024.acl-srw.49},
}

@misc{qwen38,
  author = {{Qwen Team}},
  title = {{Qwen3.8-Max}: A New Bar for Coding and Cowork},
  year = {2026},
  month = aug,
  url = {https://qwen.ai/blog?id=qwen3.8},
  note = {Model weights: Qwen3.8-27B, https://huggingface.co/Qwen/Qwen3.8-27B},
}

@article{willow2026bam,
  author = {Willow, Soohaeng Yoo and Park, Tae Hyeon and Sim, Gi Beom and Moon, Sung Wook and Min, Seung Kyu and Seo, Sangjae and Kim, Jaewook and Yang, D. ChangMo and Kim, Hyun Woo and Lee, Juho and Myung, Chang Woo},
  title = {Bayesian equivariant interatomic potential with iterative restratification of many-body message passing},
  journal = {npj Comput. Mater.},
  year = {2026},
  doi = {10.1038/s41524-026-02258-9},
}

@misc{llamacpp,
  author = {Gerganov, Georgi and {contributors}},
  title = {llama.cpp: {LLM} inference in {C/C++}},
  year = {2026},
  url = {https://github.com/ggml-org/llama.cpp},
}

@inproceedings{yang2024opro,
  author = {Yang, Chengrun and Wang, Xuezhi and Lu, Yifeng and Liu, Hanxiao and Le, Quoc V. and Zhou, Denny and Chen, Xinyun},
  title = {Large Language Models as Optimizers},
  booktitle = {The Twelfth International Conference on Learning Representations (ICLR)},
  year = {2024},
  note = {arXiv:2309.03409},
  url = {https://arxiv.org/abs/2309.03409}
}
@inproceedings{wei2022cot,
  author = {Wei, Jason and Wang, Xuezhi and Schuurmans, Dale and Bosma, Maarten and Ichter, Brian and Xia, Fei and Chi, Ed and Le, Quoc and Zhou, Denny},
  title = {Chain-of-Thought Prompting Elicits Reasoning in Large Language Models},
  booktitle = {Advances in Neural Information Processing Systems 35 (NeurIPS 2022)},
  year = {2022},
  note = {arXiv:2201.11903},
  url = {https://arxiv.org/abs/2201.11903}
}
@inproceedings{kojima2022zeroshot,
  author = {Kojima, Takeshi and Gu, Shixiang Shane and Reid, Machel and Matsuo, Yutaka and Iwasawa, Yusuke},
  title = {Large Language Models are Zero-Shot Reasoners},
  booktitle = {Advances in Neural Information Processing Systems 35 (NeurIPS 2022)},
  year = {2022},
  note = {arXiv:2205.11916},
  url = {https://arxiv.org/abs/2205.11916}
}
@misc{nye2021scratchpads,
  author = {Nye, Maxwell and Andreassen, Anders Johan and Gur-Ari, Guy and Michalewski, Henryk and Austin, Jacob and Bieber, David and Dohan, David and Lewkowycz, Aitor and Bosma, Maarten and Luan, David and Sutton, Charles and Odena, Augustus},
  title = {Show Your Work: Scratchpads for Intermediate Computation with Language Models},
  year = {2021},
  eprint = {2112.00114},
  archivePrefix = {arXiv},
  url = {https://arxiv.org/abs/2112.00114}
}
@inproceedings{vacareanu2024words,
  author = {Vacareanu, Robert and Negru, Vlad-Andrei and Suciu, Vasile and Surdeanu, Mihai},
  title = {From Words to Numbers: Your Large Language Model Is Secretly A Capable Regressor When Given In-Context Examples},
  booktitle = {First Conference on Language Modeling (COLM)},
  year = {2024},
  note = {arXiv:2404.07544},
  url = {https://arxiv.org/abs/2404.07544}
}
@misc{anthropic2026sonnet55,
  author = {{Anthropic}},
  title = {Claude Sonnet 5.5 (model documentation)},
  howpublished = {Claude Platform Docs},
  url = {https://platform.claude.com/docs/en/models/sonnet-5-5/overview},
  note = {Accessed 2026-10-15}
}
@misc{anthropic2026claudecode,
  author = {{Anthropic}},
  title = {Run Claude Code programmatically},
  howpublished = {Claude Code Docs},
  url = {https://code.claude.com/docs/en/headless},
  note = {Accessed 2026-10-15}
}
@misc{google2026gemini38flash,
  author = {{Google}},
  title = {Gemini 3.8 Flash},
  howpublished = {Gemini API documentation, Google AI for Developers},
  url = {https://ai.google.dev/gemini-api/docs/models/gemini-3.8-flash},
  note = {Accessed 2026-10-15}
}
@misc{google2026antigravitycli,
  author = {{Google}},
  title = {Headless mode, Antigravity CLI},
  howpublished = {Google Antigravity Docs},
  url = {https://antigravity.google/docs/cli/headless/},
  note = {Accessed 2026-10-15}
}
@article{luo2024cf4,
  author = {Luo, Tao and Zhang, Hang and Chen, Yingkang and Chen, Shanyong and Pan, Yang and Liu, Kang and Fu, Junwei and Chai, Liyuan and Lin, Zhang and Coote, Michelle L. and Liu, Min},
  title = {Unveiling Tetrafluoromethane Decomposition over Alumina Catalysts},
  journal = {J. Am. Chem. Soc.},
  volume = {146},
  pages = {35057--35063},
  year = {2024},
  doi = {10.1021/jacs.4c06154}
}
@article{digne2002alumina,
  author = {Digne, M. and Sautet, P. and Raybaud, P. and Euzen, P. and Toulhoat, H.},
  title = {Hydroxyl Groups on $\gamma$-Alumina Surfaces: A {DFT} Study},
  journal = {J. Catal.},
  volume = {211},
  pages = {1--5},
  year = {2002},
  doi = {10.1006/jcat.2002.3741}
}
@article{digne2004alumina,
  author = {Digne, M. and Sautet, P. and Raybaud, P. and Euzen, P. and Toulhoat, H.},
  title = {Use of {DFT} to Achieve a Rational Understanding of Acid--Basic Properties of $\gamma$-Alumina Surfaces},
  journal = {J. Catal.},
  volume = {226},
  pages = {54--68},
  year = {2004},
  doi = {10.1016/j.jcat.2004.04.020}
}
@misc{sriram2025odac25,
  author = {Sriram, Anuroop and Brabson, Logan M. and Yu, Xiaohan and Choi, Sihoon and Abdelmaqsoud, Kareem and Moubarak, Elias and de Haan, Pim and L{\"o}we, Sindy and Brehmer, Johann and Kitchin, John R. and Welling, Max and Zitnick, C. Lawrence and Ulissi, Zachary and Medford, Andrew J. and Sholl, David S.},
  title = {The Open DAC 2025 Dataset for Sorbent Discovery in Direct Air Capture},
  year = {2025},
  note = {arXiv:2508.03162},
  eprint = {2508.03162},
  archivePrefix = {arXiv},
  doi = {10.48550/arXiv.2508.03162},
}


@article{tersoff1989multicomponent, author = {Tersoff, J.}, title = {Modeling solid-state chemistry: Interatomic potentials for multicomponent systems}, journal = {Phys. Rev. B}, year = {1989}, volume = {39}, pages = {5566--5568}, doi = {10.1103/PhysRevB.39.5566},} @article{behler2007hdnn,
  author = {Behler, J{\"o}rg and Parrinello, Michele},
  title = {Generalized Neural-Network Representation of High-Dimensional Potential-Energy Surfaces},
  journal = {Phys. Rev. Lett.},
  year = {2007},
  volume = {98},
  pages = {146401},
  doi = {10.1103/PhysRevLett.98.146401},
}

@article{bartok2010gap,
  author = {Bart{\'o}k, Albert P. and Payne, Mike C. and Kondor, Risi and Cs{\'a}nyi, G{\'a}bor},
  title = {Gaussian Approximation Potentials: The Accuracy of Quantum Mechanics, without the Electrons},
  journal = {Phys. Rev. Lett.},
  year = {2010},
  volume = {104},
  pages = {136403},
  doi = {10.1103/PhysRevLett.104.136403},
}


@article{jones1924fields,
  author = {Jones, J. E.},
  title = {On the Determination of Molecular Fields. I. From the Variation of the Viscosity of a Gas with Temperature},
  journal = {Proc. R. Soc. Lond. A},
  year = {1924},
  volume = {106},
  pages = {441--462},
  doi = {10.1098/rspa.1924.0081},
}

@article{morse1929,
  author = {Morse, Philip M.},
  title = {Diatomic Molecules According to the Wave Mechanics. II. Vibrational Levels},
  journal = {Phys. Rev.},
  year = {1929},
  volume = {34},
  pages = {57--64},
  doi = {10.1103/PhysRev.34.57},
}

@misc{zang2026dac,
  author = {Zang, Yanmei and Park, Hyun Gyu and Sim, Gi Beom and Park, Tae Hyeon and Lee, Ho Jin and Zou, Xiaorong and Yang, D. ChangMo and Willow, Soohaeng Yoo and Kim, Hye Jung and Myung, Chang Woo},
  title = {Accelerated Discovery of Nitrogen-Coordinated Dual-Atom Hydrogen Evolution Reaction Electrocatalysts via Machine Learning Potentials},
  year = {2026},
  note = {arXiv:2605.29821},
  eprint = {2605.29821},
  archivePrefix = {arXiv},
  doi = {10.48550/arXiv.2605.29821},
}

@article{ha2021tuning,
  author = {Ha, Miran and Kim, Dong Yeon and Umer, Muhammad and Gladkikh, Vladislav and Myung, Chang Woo and Kim, Kwang S.},
  title = {Tuning metal single atoms embedded in {N$_x$C$_y$} moieties toward high-performance electrocatalysis},
  journal = {Energy Environ. Sci.},
  year = {2021},
  volume = {14},
  pages = {3455--3468},
  doi = {10.1039/D1EE00154J},
}

@article{hinuma2008naxcoo2,
  author = {Hinuma, Yoyo and Meng, Ying S. and Ceder, Gerbrand},
  title = {Temperature-concentration phase diagram of {P2}-{Na$_x$CoO$_2$} from first-principles calculations},
  journal = {Phys. Rev. B},
  year = {2008},
  volume = {77},
  pages = {224111},
  doi = {10.1103/PhysRevB.77.224111},
}

@article{zheng2017p2mixing,
  author = {Zheng, Chen and Radhakrishnan, Balachandran and Chu, Iek-Heng and Wang, Zhenbin and Ong, Shyue Ping},
  title = {Effects of Transition-Metal Mixing on Na Ordering and Kinetics in Layered {P2} Oxides},
  journal = {Phys. Rev. Appl.},
  year = {2017},
  volume = {7},
  pages = {064003},
  doi = {10.1103/PhysRevApplied.7.064003},
}

@article{brown1985bvs,
  author = {Brown, I. D. and Altermatt, D.},
  title = {Bond-valence parameters obtained from a systematic analysis of the Inorganic Crystal Structure Database},
  journal = {Acta Crystallogr. B},
  year = {1985},
  volume = {41},
  pages = {244--247},
  doi = {10.1107/S0108768185002063},
}

\end{filecontents*}

\documentclass[journal=jacsat,manuscript=article]{achemso}

\author{Hyun Gyu Park}
\affiliation{Department of Energy Science, Sungkyunkwan University, Suwon, Republic of Korea}
\author{Won Woong Choi}
\affiliation{Department of Energy Science, Sungkyunkwan University, Suwon, Republic of Korea}
\author{Wonho Lee}
\affiliation{Department of Energy, Sungkyunkwan University, Suwon, Republic of Korea}
\author{Soohaeng Yoo Willow}
\affiliation{Department of Energy Science, Sungkyunkwan University, Suwon, Republic of Korea}
\author{Sangjae Seo}
\affiliation{Department of Supercomputing Acceleration Research, Korea Institute of Science and Technology Information, Daejeon 34141, Republic of Korea}
\author{Juho Lee}
\email{juholee@kaist.ac.kr}
\affiliation{Kim Jaechul Graduate School of AI, KAIST, Daejeon, Republic of Korea}
\author{Chang Woo Myung}
\email{cwmyung@skku.edu}
\affiliation{Department of Energy Science, Sungkyunkwan University, Suwon, Republic of Korea}
\alsoaffiliation{Department of Energy, Sungkyunkwan University, Suwon, Republic of Korea}

\title{Language-Model-Written Physical Corrections for a Pretrained Interatomic Potential}

\usepackage{booktabs}
\usepackage{graphicx}
\usepackage{amsmath,amssymb}
\usepackage{tabularx}
\usepackage{tikz}
\usetikzlibrary{arrows.meta}
\usepackage{fvextra}
\usepackage{placeins}

% Battery and bulk test sets (Table "Test sets" and its results): \bulktrue shows them, \bulkfalse hides them.
\newif\ifbulk
\bulkfalse

\begin{document}

\begin{abstract}
\ifbulk
Pretrained interatomic potentials can carry a large energy offset on a new system, and per element constants that remove this offset leave the forces unchanged. We ask whether a frozen language model can write a correction that improves both energies and forces, with every number chosen by the model rather than by a fitting routine. The neural network energy of BAM is kept frozen, and the language model receives 20 labeled examples and writes per element constants together with physically named terms: Born--Mayer repulsion, Morse bonds, embedded atom and bond angle terms, with parameter ranges set from the example structures. Candidates from ten calls are scored on the examples; only candidates with lower energy and force errors than BAM are kept, and the one with the lowest force error is applied unchanged to the queries. Across eight test sets from five materials (Al--Mg--Si alloys, $\mathrm{Li_6PS_5Cl}$, a Ba--Sc--W perovskite proton conductor, $\mathrm{B_{32}N_{32}}$ and Na--Mn--O), each with 50 queries, the energy mean absolute error falls by 68--99\%, for example from 285.09 to 2.96 meV/atom, and the force mean absolute error falls by 3.9--18.3\%, with paired 95\% bootstrap intervals that exclude zero in every set. Every selected correction uses only Born--Mayer and Morse terms. Energy differences between structures of the same formula change by at most 2.1 meV/atom, so most of the energy gain is a correction of the per element reference. On identical random sets of 20, 10 and 5 examples and with four language models, full fine tuning of BAM gives lower energy and force errors for $\mathrm{Li_6PS_5Cl}$ and lower force errors for the proton conductor, fine tuning degrades the Al alloys while the LM written corrections stay at the BAM level, and the comparison is mixed for $\mathrm{B_{32}N_{32}}$ and Na--Mn--O. We also test the method on the hydrogen adsorption energy of dual atom catalysts M$_2$-N$_x$-C (M = Co, Fe, Cr, Ir, Pt), in which per element constants cancel. With the adsorption energy stated as the target and a check on short range forces, the five run Claude mean lowers the adsorption energy MAE on unseen N coordination motifs from 0.428~eV to 0.220--0.247~eV with 5 to 20 examples. Fine tuning on the same examples gives lower force errors but adsorption energy errors of 0.406--0.509~eV. For CO$_2$ adsorption in 33 metal organic frameworks, the same approach with an added dispersion term lowers the adsorption energy error at an unseen site of each framework from 0.260 to 0.153~eV with 20 examples, whereas fine tuning raises it to 0.337~eV, and a correction written for one framework does not transfer to others. For CF$_4$ on three Al$_2$O$_3$ surfaces, the same approach without the dispersion term lowers the adsorption energy error at seven unseen sites from 0.228 to 0.082~eV, with an F--Al repulsion on every surface, and five runs each with Claude, Gemini and GPT stay between 0.071 and 0.168~eV, whereas fine tuning reaches 0.183~eV. For P2 Na$_x$MnO$_2$, where the per element constants cancel in the formation energy, a version with 35 calls lowered the formation energy error from 59.7 to 20.9--31.4~meV per formula unit in four of five runs with Claude Sonnet 5.5 but made the ordering of structures at the same composition worse, and 17 of 20 runs of four language models limited to ten calls kept the BAM prediction. Apart from the cathode test, the main results come from one language model with one run per set, the cathode, catalyst, framework and alumina variants were designed after the first results for each system had been seen, and the results do not show a more accurate potential energy surface.
\else
Pretrained interatomic potentials can exhibit system specific energy prediction errors when applied to new materials or chemical environments. Fine tuning is commonly used to reduce these errors, but requires updating the pretrained neural network weights. Here, four frozen language models (Claude Sonnet 5.5, Gemini 3.8 Flash, GPT-5.6 Luna and Qwen3.8-27B) write explicit physical corrections to BAM-MP-core (BAM) without updating its neural network weights. They select terms from a prescribed physical library and write all adjustable parameters from a small set of labeled examples. Candidates are scored on those examples and applied without subsequent coefficient fitting. With 21 examples, including the Na metal reference, the P2 Na$_x$MnO$_2$ cathode formation energy mean absolute error (MAE) is 31.68--54.22~meV per formula unit (BAM: 59.66; fine tuning: 39.13). With 20 examples per metal, the metal averaged hydrogen adsorption energy MAE for dual atom catalysts based on the five evaluated metals is 0.220--0.342~eV (BAM: 0.428; fine tuning: 0.406). With 20 examples per framework, the CO$_2$ adsorption energy MAE in 33 metal organic frameworks is 0.162--0.215~eV (BAM: 0.260; fine tuning: 0.337); in 11 of these frameworks the test CO$_2$ lies at an example site, and on the other 22 it is 0.194--0.228~eV (BAM: 0.261; fine tuning: 0.315). With 6--10 examples per surface, the CF$_4$ adsorption energy MAE across seven sites on three alumina surfaces is 0.086--0.182~eV (BAM: 0.228; fine tuning: 0.183). These results demonstrate that language model written physical corrections improve energy difference predictions relative to uncorrected BAM for the evaluated downstream tasks, without neural network retraining.
\fi
\end{abstract}

\section{Introduction}

Atomistic simulations require accurate energies and forces at a computational cost that permits large systems and long trajectories. Electronic structure methods such as density functional theory (DFT) provide a first principles description, but their cost limits these scales. Classical interatomic potentials, including the Lennard--Jones and Morse potentials, reduce this cost through prescribed functional forms with parameters fitted to data.\cite{jones1924fields,morse1929} Their limited flexibility can constrain predictive accuracy.\cite{batzner2022} Machine learned potentials use more flexible neural network or kernel representations trained on electronic structure energies and forces. Early neural network potentials and Gaussian approximation potentials demonstrated that electronic structure predictions could be reproduced at a much lower evaluation cost.\cite{behler2007hdnn,bartok2010gap}

Early neural network and Gaussian approximation potentials were trained for specific chemical systems. Applying them to other chemistries often required new electronic structure data and training.\cite{chen2022m3gnet} Pretraining on large datasets spanning many elements and compositions has since produced potentials such as M3GNet, CHGNet, MACE-MP-0 and BAM-MP-core.\cite{chen2022m3gnet,deng2023chgnet,batatia2025macemp0,willow2026bam} Broad coverage does not, however, ensure accurate energies and forces for a particular target system.\cite{deng2025softening,focassio2025surfaces} Fine tuning addresses these errors by updating pretrained weights with target system data rather than training from scratch.\cite{radova2025frozen} Such updates can degrade predictions for previously learned systems, an effect known as catastrophic forgetting.\cite{kim2025reewc}

An alternative to fine tuning is to construct an explicit correction from a small set of labeled structures while keeping the pretrained neural network weights fixed. Language models have been used to propose numerical solutions\cite{vacareanu2024words,yang2024opro} and to generate equations and programs.\cite{merler2024icsr,shojaee2025llmsr,romera2024funsearch} These capabilities motivate their use in writing complete physical corrections, including both the interaction terms and their numerical coefficients. The relevant test is whether complete expressions, with their coefficients used as written, improve application specific energy differences rather than merely align total energies.

Our main contributions are summarized as follows:
\begin{enumerate}
\item We introduce a method for task specific adaptation of the pretrained BAM-MP-core potential\cite{willow2026bam} through explicit physical corrections written by frozen language models, without updating the pretrained neural network weights (Fig.~\ref{fig:schematic}).
\item We develop correction procedures tailored to each application, in which language models select terms from a prescribed physical library and write all numerical coefficients from labeled examples and prediction error summaries, without external coefficient fitting.
\item \ifbulk We evaluate absolute energies and forces on eight test sets from five materials, as well as the formation energy, convex hull and voltage of a sodium ion cathode and three adsorption energies.\else We benchmark four language models on four downstream applications: formation energies, phase stability and voltage in a P2 Na$_x$MnO$_2$ cathode; hydrogen adsorption on dual atom catalysts; CO$_2$ adsorption in metal organic frameworks; and CF$_4$ adsorption on alumina surfaces.\fi{} We compare with unmodified BAM and potentials fine tuned on the same labeled examples, repeating the complete correction procedure five times per language model using fixed example and evaluation sets to obtain more reliable performance estimates.
\end{enumerate}

Rather than reoptimizing the pretrained neural network weights, we formulate task specific adaptation as the construction of an explicit physical correction. In this study, we therefore develop and evaluate a method in which frozen language models write the complete correction, including its numerical coefficients, directly from labeled examples without external fitting.

\begin{figure}[p]
\centering
\includegraphics[width=\linewidth]{bam-icl-schematic-blender-20261023.pdf}
\caption{Workflow for language model written physical corrections to a frozen interatomic potential. (a) System specific labeled examples provide elemental atom counts, DFT energies and forces, neural network outputs, and residual and geometric summaries. (b) A frozen language model writes complete candidate expressions, including all numerical coefficients, using prescribed physical forms and parameter bounds. Candidates are scored and refined on the examples according to the target energy objective and acceptance criteria of each application. The best passing expression is applied without parameter refitting, and unmodified BAM-MP-core is used if none passes. (c) The selected expression replaces the pre existing external elemental reference sum and is evaluated at the query geometry with fixed neural network weights, giving an energy correction and its negative coordinate gradient as a force correction. Labels of the evaluation structures are not used during generation and selection and are read only after the selected correction is saved. The atomic renderings are illustrative.}
\label{fig:schematic}
\end{figure}

\section{Methods}

\subsection{Energy assembly and physical corrections}

BAM denotes the unmodified BAM-MP-core model. For atomic coordinates $\mathbf R$ and atomic numbers $\mathbf Z$, its total energy is
\begin{equation}
\begin{aligned}
E_{\mathrm{BAM}}(\mathbf R,\mathbf Z)&=E_{\mathrm{NN}}(\mathbf R,\mathbf Z)+R_{\mathrm{old}}(\mathbf Z),\\
R_{\mathrm{old}}(\mathbf Z)&=\sum_i\left[\overline e(Z_i)+c_{\mathrm{old}}(Z_i)\right],
\end{aligned}
\label{eq:bam}
\end{equation}
Here, $E_{\mathrm{NN}}=\sum_i e_i^{\mathrm{NN}}$ is the sum of atomic energy contributions obtained from BAM's message passing features through its pretrained nonlinear readout layers. The fitted elemental reference energies $\overline e(Z)$ and the element wise energy shift corrections $c_{\mathrm{old}}(Z)$ are stored in the pretrained BAM checkpoint. Since $R_{\mathrm{old}}$ depends only on composition, it does not contribute to forces, giving $\mathbf F_{\mathrm{BAM}}=-\nabla_{\mathbf R}E_{\mathrm{NN}}$.

A frozen language model (LM) writes an explicit physical correction from labeled examples, called the LM written correction below. The correction has the form
\begin{equation}
f(\mathbf R,\mathbf Z)=\sum_Z n_Zc_Z+\sum_{t=1}^{T}\phi_t(\mathbf R,\mathbf Z),\qquad 1\leq T\leq8,
\label{eq:f}
\end{equation}
where $n_Z$ is the number of atoms of element $Z$, $c_Z$ is an elemental energy constant, and $\phi_t$ is a coordinate dependent physical term. Replacing $R_{\mathrm{old}}$ with $f$ gives the corrected energy and forces,
\begin{equation}
\begin{aligned}
E_{\mathrm{corr}}&=E_{\mathrm{NN}}+f,\\
\mathbf F_{\mathrm{corr}}&=-\nabla_{\mathbf R}E_{\mathrm{corr}}=\mathbf F_{\mathrm{BAM}}-\nabla_{\mathbf R}f.
\end{aligned}
\label{eq:new}
\end{equation}
For each prediction, $f$ is evaluated at the atomic coordinates and atomic numbers of the structure of interest. The elemental constants change energies but not forces. Force corrections arise only from the coordinate dependent terms.

The common physical library comprises Born--Mayer, Morse, embedded atom and bond angle forms. These are
\begin{align}
\phi^{\mathrm{BM}}_{XY}&=\sum_{i<j\in XY}B\,e^{-(r_{ij}-r_{1,XY})/\rho}\,S(r_{ij}), \label{eq:bm}\\
\phi^{\mathrm{Morse}}_{XY}&=\sum_{i<j\in XY}D\left[\left(1-e^{-a(r_{ij}-r_0)}\right)^2-1\right]S(r_{ij}), \label{eq:morse}\\
\phi^{\mathrm{EAM}}_X&=-A\sum_{i\in X}\sqrt{\rho_i},\qquad \rho_i=\sum_j e^{-\beta(r_{ij}-r_{1,X})}S(r_{ij}), \label{eq:eam}\\
\phi^{\mathrm{angle}}_X&=\lambda\sum_{i\in X}\sum_{j<k}\left(\cos\theta_{jik}-\cos\theta_0\right)^2S_a(r_{ij})S_a(r_{ik}). \label{eq:angle}
\end{align}
For Born--Mayer and Morse, $XY$ specifies the element pair used in the pair sum. The embedded atom contribution is summed over atoms of element $X$. In the bond angle term, atoms of element $X$ are the angle centres.

The spatial range of the four common terms is set by a neighbour list, two reference distances and switching functions. All four terms use a neighbour list with a 5~\AA\ cutoff, which includes periodic images along periodic directions. The two reference distances are computed from the example structures alone, without query structures or DFT labels. For an element $X$, $r_{1,X}$ is the median, over all $X$ atoms in the examples, of the distance to the nearest neighbour of any element. For an element pair $XY$, $r_{1,XY}$ is the median of one pooled list that contains the distance from each $X$ atom to its nearest $Y$ atom and from each $Y$ atom to its nearest $X$ atom. For $X=Y$, the list contains the distance from each $X$ atom to its nearest $X$ neighbour. Only neighbours within the cutoff are considered. The distance $r_{1,XY}$ is the reference distance in Eq.~\eqref{eq:bm} and scales the allowed range of $r_0$ in Eq.~\eqref{eq:morse}, and $r_{1,X}$ is the reference distance in Eq.~\eqref{eq:eam} and scales the range of the angle term in Eq.~\eqref{eq:angle}.

 For the Born--Mayer and Morse terms and the embedded atom density, each neighbour contribution is multiplied by a cosine switching function $S(r)$. This smoothly removes the contribution near the cutoff rather than truncating it abruptly.\cite{tersoff1989multicomponent,behler2007hdnn} The switching function $S(r)$ equals 1 for $r\le4.5$~\AA\ and 0 for $r\ge5$~\AA. Between 4.5 and 5~\AA, it decreases smoothly from 1 to 0 according to a cosine function. These distances are fixed implementation settings.
\[
S(r)=\begin{cases}1, & r\le4.5\,\text{\AA},\\[2pt]\frac12\!\left[1+\cos\!\left(\pi\frac{r-4.5\,\text{\AA}}{0.5\,\text{\AA}}\right)\right], & 4.5\,\text{\AA}<r<5\,\text{\AA},\\[2pt]0, & r\ge5\,\text{\AA}.\end{cases}
\]
The bond angle term uses an element specific range defined relative to $r_{1,X}$. Only neighbours satisfying $r<1.35r_{1,X}$ are retained from the common neighbour list. Its switching function $S_a(r)$ has the same cosine form. It equals 1 for $r\le1.15r_{1,X}$, decreases smoothly to 0 between $1.15r_{1,X}$ and $1.35r_{1,X}$, and remains zero beyond that range. These cutoff distances and scaling factors are fixed implementation settings.

The language model selects the term forms and their element pair or central element targets and writes all numerical coefficients. Each of the four forms may occur at most once for a given target. Parameters are restricted to prescribed ranges, including $B,D,A,\lambda\in[-1,1]$~eV and $r_0\in[0.8,1.2]r_{1,XY}$, and the complete ranges are given in Table~S1. Details of the procedures and additional correction terms used for each application are provided in Section~S2 of the Supporting Information.

Candidate expressions are parsed, checked against the prescribed format and parameter ranges, and evaluated with their coefficients as written. No external regression or numerical optimizer fits or adjusts these coefficients. Force corrections are calculated by automatic differentiation in double precision. Neither the pretrained interatomic potential weights nor the language model weights are updated.

\ifbulk
\subsection{Datasets}

Table~\ref{tab:splits} lists eight test sets from five materials. Each has 50 evaluation queries, three disjoint development queries and 20 examples per query. Queries that share examples share one example context, and the language model writes one correction per context.

The Al queries are the 50-frame panel of the earlier protocol, drawn with the NumPy generator \texttt{default\_rng(12345)} from a 564-frame eligible test bundle. The structures contain 2--33 atoms of Mg, Al and Si, and examples come only from the Al training data, retrieved by composition. The 50 queries use 25 example contexts.

The $\mathrm{Li_6PS_5Cl}$ (LPSCl) structures contain 52 atoms, $\mathrm{Li_{24}P_4S_{20}Cl_4}$. For interpolation the examples come from 600--800 and 1000--1200~K and the queries from 900~K. For extrapolation the examples come from 600--800~K and the queries from 1000--1200~K. The BSW structures contain 138--162 atoms of H, O, Sc, Ba and W. For temperature transfer the examples are labeled 550~$^{\circ}$C and the queries 600~$^{\circ}$C (seven contexts). For composition transfer the queries carry the nominal W label 0.15 and the examples the other labels. All 20 retrieved examples have one 153-atom composition, whereas the queries have 150 atoms, so every query lies outside the composition span of its examples. The $\mathrm{B_{32}N_{32}}$ (BN) set contains 64-atom structures. Its 53 test frames and 20 training frames were drawn with \texttt{default\_rng(12345)} without reading labels, and all queries share the 20 examples.

The $\mathrm{Na}_x\mathrm{MnO_2}$ (NMO) structures contain 57--118 atoms of Na, Mn and O in P2, P2$'$ and OP4 stackings with $x$ between 0 and 1. Training and test frames come from disjoint sets of calculations (587 and 73 runs with 28\,217 and 3804 frames). The 53 test frames were drawn with \texttt{default\_rng(12345)} without reading labels. For each query, 20 training frames are retrieved by the $L_1$ distance between element fractions, with at most two frames per run, at least three runs and every element of the query present. In the phase restricted set the examples must also share the stacking of the query. The two NMO sets use 16 and 19 contexts, and the LPSCl, BN and BSW composition sets use one.

\begin{table}[t]
\caption{Test sets. Each uses 50 evaluation queries, three development queries and 20 examples per query. Contexts are distinct example sets.}
\label{tab:splits}
\small
\begin{tabularx}{\textwidth}{l>{\raggedright\arraybackslash}X>{\raggedright\arraybackslash}Xr}
\toprule
Test set & Example pool & Queries & Contexts \\
\midrule
Al alloys & Al training data & 50 of 564 eligible Al test frames & 25 \\
LPSCl interpolation & 600--800 and 1000--1200~K, 2093 frames & 900~K & 1 \\
LPSCl extrapolation & 600--800~K, 1065 frames & 1000--1200~K & 1 \\
BSW temperature & 550~$^{\circ}$C, 200 frames & 600~$^{\circ}$C & 7 \\
BSW composition & Nominal W label $\ne0.15$, 342 frames & W label 0.15 & 1 \\
BN & 20 of 40\,000 training frames & 50 of 5000 test frames & 1 \\
NMO composition & 28\,217 training frames, any stacking & 50 of 3804 test frames & 16 \\
NMO composition and phase & 28\,217 training frames, same stacking as the query & 50 of 3804 test frames & 19 \\
\bottomrule
\end{tabularx}
\end{table}

\textit{Limitations of data provenance.} The original Al electronic structure settings are not established by the stored frame records. The BSW and BN energy and force origins and their common electronic structure settings remain unverified, so their labels are treated as supplied labels, not as an independently verified DFT benchmark. The LPSCl split manifest documents temperature and subsampling but does not establish the electronic structure settings, and the NMO frames are read from calculation outputs whose settings are not reported here. The imported historical bundles also have an unresolved producer-code hash discrepancy. Their artifact hashes and current import contracts are checked, but this is not retrospective verification of the historical producer implementation.

\fi

\subsection{Correction generation and selection}

Without updating their weights, we evaluate four language models, Claude Sonnet 5.5,\cite{anthropic2026sonnet55,anthropic2026claudecode} Gemini 3.8 Flash,\cite{google2026gemini38flash} GPT-5.6 Luna and Qwen3.8-27B.\cite{qwen38,llamacpp} Claude uses reasoning, Gemini uses the low effort variant, and GPT uses medium reasoning effort. Qwen is run locally with 4 bit (Q4\_K\_M) quantization and reasoning disabled. No tools are available to the models, so they cannot run code to fit the coefficients. Execution details and response handling are given in Section~S1.2 of the Supporting Information.

In the base protocol, all models receive the same initial instructions and data from 20 examples, comprising composition, structural summaries and the energy and force residuals of BAM relative to DFT. Neither query geometries nor their labels are provided to the models. Each model writes a correction function with explicit numerical coefficients, which are used without subsequent fitting. Each example set, called a context below, receives five initial calls followed by five refinement calls, with up to five candidates per call. Refinement prompts add the five best previous candidates and their errors on the examples. These inputs depend on each model’s responses, but calls carry no conversation history. A candidate is eligible only if both its energy and force errors on the examples are lower than those of BAM. Among eligible candidates, the one with the lowest force error is selected. If none is eligible, BAM is used. Generation and selection use the same examples, without access to query targets. The full prompt and descriptor definitions are provided in Section~S1, and changes to the inputs and selection criteria for each application are described in Section~S2.

Each application adds its target quantity to the prompt and selects candidates by their example error in that quantity. For the sodium ion cathode, the prompt includes the DFT and BAM formation energies of the examples, and the selection rule targets the hull, as described in the cathode subsection below. For the three adsorption applications, namely the dual atom catalysts, the MOFs and CF$_4$ on alumina, the prompt includes the DFT and BAM adsorption energies of the examples. In these applications, a candidate is eligible only if its energy, force and adsorption energy errors on the examples are all lower than those of BAM and it passes a range check. In this check, no term may push or pull two atoms placed alone with a force above 3~eV/\AA\ at any distance from close contact (0.9 times the sum of their covalent radii) to 5~\AA\ (Section~S2.2). Among the candidates that pass, the one with the lowest example adsorption energy error is selected. If none passes, BAM is retained. The adsorption comparisons use five runs per language model and ten calls per example set, with fixed examples and evaluation structures.

\ifbulk
\subsection{Evaluation}

For the $Q$ query structures, we calculate the mean absolute energy and force errors as
\begin{align}
\operatorname{MAE}_E&=\frac1Q\sum_q\frac{|E_q-E_q^\star|}{N_q},\\
\operatorname{MAE}_{F}&=\frac1Q\sum_q\frac1{N_q}\sum_i\|\mathbf F_{qi}-\mathbf F_{qi}^\star\|_2 ,
\end{align}
Here, $E_q$ is the predicted total energy for query $q$, and $E_q^\star$ is its DFT reference energy. $N_q$ is the number of atoms in the query. The energy MAE is the average absolute energy error per atom across the queries and is reported in meV/atom. For atom $i$ in query $q$, $\mathbf F_{qi}$ and $\mathbf F_{qi}^\star$ are the predicted and DFT reference force vectors. The force MAE is reported in eV/\AA.
\fi

\ifbulk
To measure errors in relative energies, we compare pairs of queries with the same proportions of elements. For each query, we calculate the per atom energy error by subtracting its DFT reference energy from its predicted total energy and dividing by the number of atoms. For each pair, we take the absolute difference between the two energy errors per atom. We average these pair errors and count each pair once. Let $\mathcal{P}$ denote the set of these pairs. In each pair, $q$ and $q'$ are the indices of two distinct query structures. Then
\begin{equation}
\operatorname{MAE}_{\Delta E}=\frac{1}{|\mathcal{P}|}\sum_{(q,q')\in\mathcal{P}}\left|\frac{E_q-E_q^\star}{N_q}-\frac{E_{q'}-E_{q'}^\star}{N_{q'}}\right|.
\end{equation}
Thus, elemental ratios with more query pairs contribute more to the result. 

To estimate sampling uncertainty in the energy and force improvements over BAM, we draw 10\,000 samples of query frames with replacement, using seed 12345. In each sample, we calculate the errors of BAM with and without the LM written correction on the same frames, then define the paired improvement as the error without the correction minus the error with it. We report 95\% intervals for these improvements, and positive values indicate lower errors after correction. Queries that share a context also share the same LM written correction. Because the samples vary query frames but hold model calls and example sets fixed, these intervals do not capture variation across contexts, repeated runs, or new example sets.\ifbulk Each test set was run once.\fi
\fi

\ifbulk We reused the query sets from the earlier protocol, so they are not new test data.\fi{} 

\ifbulk
\subsection{Comparison with fine tuning}

To compare the LM written correction with fine tuning on identical data, 20 labeled frames were drawn at random for each material with \texttt{default\_rng(0)}, and nested subsets of 10 and 5 frames were taken from them. None of these frames coincides with an evaluation query, and both methods received the same frames. One example set per material and size serves all queries of that material, instead of the per query retrieval of Table~\ref{tab:splits}, so the language model errors in this comparison differ from those in Table~\ref{tab:main}. Four language models followed the procedure described above, with five initial and five refinement calls per example set. Claude Sonnet 5.5 used the settings of the main runs. Gemini 3.8 Flash\cite{google2026gemini38flash} used the low effort variant and GPT-5.6 Luna medium reasoning effort, both through their command line interfaces with tools disabled. Qwen3.8-27B\cite{qwen38} was run locally as a 4-bit quantized model with llama.cpp\cite{llamacpp} and reasoning off. Its sampling seed is derived from the example set and the call, so a repeated run on the same examples would return the same answers. Each set was run once. All four received identical initial prompts for a given example set. Fine tuning updated all weights of BAM, starting from the pretrained checkpoint, with a mean absolute error loss that weights energies and forces equally, batch size 2 and 200 epochs. The learning rate started at $10^{-4}$ and was lowered to 70\% of its current value whenever the validation loss had not improved for 20 consecutive epochs. Of the 20, 10 and 5 frames, 16, 8 and 4 were used for training and 4, 2 and 1 for selecting the epoch with the lowest validation loss. The per element reference energies of the fine tuned model were refitted on all frames of the set. The fine tuning settings were fixed before the runs and were not tuned. Because one LM written correction and one fine tuned model serve all queries of a test set, this comparison is evaluated on larger query sets than Table~\ref{tab:main}, namely all 564 frames of the eligible Al test bundle, all 200 frames of each LPSCl query pool, all 200 BSW frames labeled 600~$^{\circ}$C, all 58 BSW frames with the W label 0.15, and the first 1000 frames of the \texttt{default\_rng(12345)} permutation of the BN and NMO test frames. Each set contains the 50 queries of Table~\ref{tab:main}, and none of its frames is an example of this comparison. The two NMO test sets share their queries in this comparison and are reported once.

For BSW, the first random draw of 20 frames contained eight frames labeled 600~$^{\circ}$C, two of them with the W label 0.15, which belong to the query domains of the two BSW test sets. We kept its first ten frames, which are all labeled 550~$^{\circ}$C with other W labels and form the 10- and 5-frame subsets, replaced the other ten by frames drawn with \texttt{default\_rng(0)} from the 152 eligible 550~$^{\circ}$C frames with W labels other than 0.15, and reran both methods on the corrected set. No example in this comparison comes from a query domain. One further feature of the draws affects the comparison. For Al, the frames include structures containing Fe, Mn or Zn, elements that do not occur in the Mg--Al--Si queries. Each condition was run once with a single draw.

\else
\subsection{Comparison with fine tuning}

Fine tuning updates all BAM weights from the pretrained checkpoint using a mean absolute error loss with equal energy and force weights, batch size 2 and 200 epochs. It uses the same examples as the LM written correction, and a separate model is trained once for each context and each example set size. A few examples are set aside as a validation set and are not used for training. In the adsorption applications, the bare host and the isolated molecule always remain in training (Section~S2). After each epoch, the error on the validation set, called the validation loss, is computed, and the model from the epoch with the lowest validation loss is kept. The learning rate starts at $10^{-4}$ and is lowered to 70\% of its current value whenever the validation loss has not improved for 20 consecutive epochs. The fine tuned model's per element reference energies are refitted on all example structures in each set.

\fi

\subsection{Evaluated quantities}
\label{sec:application-targets}

Energies and forces are evaluated at the DFT geometries. These structures are not relaxed again with BAM, the fine tuned model or BAM with the LM written correction.\ifbulk\else{} Energy MAE is the mean over evaluation structures of $|E-E^\star|/N$, where $E$ and $E^\star$ are the predicted and DFT energies and $N$ is the number of atoms. It is reported in meV/atom. Force MAE is the norm of the force vector error averaged over the atoms of each structure and then over structures, reported in eV/\AA.\fi

Adsorption energies are
\begin{equation}
E_{\mathrm{ads}}=E(\text{host+adsorbate})-E(\text{host})-\nu E(\text{reference molecule}),
\label{eq:adsorption}
\end{equation}
where $\nu=1/2$ for H adsorption, referenced to the H$_2$ molecule, and $\nu=1$ for CO$_2$ and CF$_4$ adsorption, referenced to the isolated CO$_2$ and CF$_4$ molecules. BAM, the fine tuned model and BAM with the LM written correction each evaluate the three energies in Eq.~\eqref{eq:adsorption} with one energy function and one set of coefficients. Per element constants cancel in these adsorption energies and in the formation energies and voltages of the sodium ion cathode below, so any change from BAM in these quantities comes from the physical terms $\phi_t$.

\subsection{Sodium ion cathode}
\label{sec:hull-methods}

For P2 $\mathrm{Na}_x\mathrm{MnO_2}$ (NMO), we use the formation energy per MnO$_2$ formula unit relative to the fully desodiated and fully sodiated end members,
\begin{equation}
\Delta E_\mathrm{f}(x)=E(\mathrm{Na}_x\mathrm{MnO_2})-x\,E(\mathrm{NaMnO_2})-(1-x)\,E(\mathrm{MnO_2}),
\label{eq:emix}
\end{equation}
where the end members are the P2 structures at $x=0$ and $x=1$. Structures on the lower convex hull of $\Delta E_\mathrm{f}$ versus $x$ are the stable P2 structures of this set at 0~K, and $E_{\mathrm{hull}}$ is the energy above this hull. The voltage between adjacent vertices $x_1<x_2$ of the DFT hull is
\begin{equation}
V=-\frac{E(\mathrm{Na}_{x_2}\mathrm{MnO_2})-E(\mathrm{Na}_{x_1}\mathrm{MnO_2})-(x_2-x_1)\,E(\mathrm{Na})}{(x_2-x_1)\,e},
\label{eq:voltage}
\end{equation}
where $E(\mathrm{Na})$ is the energy per atom of Na metal. Each model uses its own structure energies and Na metal energy (Section~S2.1 of the Supporting Information).

The P2 NMO reference energies are computed with PBE$+U$ ($U=4$~eV on Mn) and the D3(BJ) dispersion correction. The 21 examples consist of 20 P2 structures, including both end members, and Na metal. The 279 evaluation structures consist of the 277 P2 structures that are not examples and the two end members, which are also examples. The evaluation set is therefore not fully independent of the examples. Each run yields one correction, which is applied to all 279 evaluation structures and to Na metal, so that the hull and the voltages are computed with one energy function.

The cathode version differs from the base protocol in the prompt, the term library, the selection rule and the number of calls. It also uses 21 examples, including Na metal, rather than 20. First, the prompt states $\Delta E_\mathrm{f}$ of Eq.~\eqref{eq:emix} as the target and gives its DFT and BAM values for the examples. It splits the BAM error of $\Delta E_\mathrm{f}$ at the example compositions into a depth, proportional to $x(1-x)$ and symmetric about $x=0.5$, and a tilt, proportional to $x(1-x)(1-2x)$ and of opposite sign on the Na poor and Na rich sides (Section~S2.1). It also reports how a small term of each type changes these components. Every number in the prompt that is computed from structures comes from the 21 examples, and the other numbers are fixed settings such as parameter ranges and thresholds. No evaluation structure or label enters the prompt.

Second, besides the four common terms, the library offers four terms aimed at Na ordering and Mn oxidation (Section~S2.1). A Gaussian shell term adds a positive or negative energy to each pair of two chosen elements, largest at a chosen neighbour distance and decaying smoothly away from it. A screened Coulomb term adds a repulsion between two chosen elements that reaches farther than the other terms and is summed up to 10~\AA. A site term counts, for each atom of a chosen element, the Mn atoms within a chosen distance, with a smooth cutoff, and adds the same chosen energy, positive or negative, for each counted Mn atom. With Na as this element, it can distinguish the two Na sites of P2 layered oxides, including Na$_x$MnO$_2$. One site shares faces and the other only edges with the transition metal octahedra, and the two sites differ in energy.\cite{hinuma2008naxcoo2,zheng2017p2mixing} A bond valence term estimates the oxidation state of each Mn atom from its Mn--O distances, using the fixed Mn(IV)--O parameters of Brown and Altermatt for every Mn atom,\cite{brown1985bvs} and adds an energy proportional to the squared deviation of this value from a chosen reference.

Third, selection targets the hull rather than the forces. On the examples, a candidate passes only if its errors in energy, $\Delta E_\mathrm{f}$ and $E_\mathrm{hull}$ are lower than those of BAM, its force error is at most 1.2 times that of BAM, it orders pairs of structures with the same $x$ as DFT does at least as often as BAM, the magnitude of the tilt of its remaining $\Delta E_\mathrm{f}$ error is at most 0.7 times that of BAM, and it passes the range check, which for the screened Coulomb term extends to 10~\AA. Among the passing candidates, the one with the lowest sum of the example $E_\mathrm{hull}$ error and the remaining tilt is selected. If no candidate passes, BAM is retained.

Fourth, each run uses 35 calls, five initial and 30 refinement calls, instead of the ten calls of the base protocol, and five runs are performed per language model. Ten calls rarely yield a candidate that passes the cathode selection rule, as shown in the cathode results and Table~S\ifbulk 4\else 2\fi. The prompt was revised after seeing results on the evaluation structures (Section~S2.1).

\subsection{Dual atom catalysts}
\label{sec:catalyst-methods}

The catalysts M$_2$-N$_x$-C are single graphene sheets that hold two metal atoms and zero to six N atoms, for M = Co, Fe, Cr, Ir and Pt, taken from the data set of Zang et al.\cite{zang2026dac} The reference energies are spin polarized RPBE with Tkatchenko and Scheffler dispersion and a Hubbard $U$ of 3.42, 3.29 and 2.79~eV on Co, Fe and Cr, taken from Ha et al.,\cite{ha2021tuning} and no $U$ on Ir and Pt. The target is the hydrogen adsorption energy $\Delta E_\mathrm{H}=E_\mathrm{ads}$ of Eq.~\eqref{eq:adsorption}, with $\nu=1/2$ and H$_2$ as reference.

Each metal has 20 N coordination motifs. Ten supply the examples, 5, 10 or 20 slab and slab+H structures, and the remaining ten were reserved for testing, of which 4N/c was excluded for every metal. Both the language models and fine tuning receive the DFT H$_2$ molecule as an additional example, because no slab+H example contains an H--H pair. BAM uses no examples. Energy and force errors were evaluated on 135 test structures: 45 bare catalyst slabs, 45 relaxed H-adsorbed slabs and 45 unrelaxed H-adsorbed structures evaluated by single-point calculations. The 45 pairs of bare and relaxed H-adsorbed slabs were used to evaluate adsorption energies. The Cr 4N/c slab+H file was a duplicate of a V calculation, and some single points were computed without $U$ (Section~S2.2). Metal selection used BAM errors including the test motifs and applies equally to all methods. For each metal and example set size, one correction is written from the examples of that metal and applied to its nine test pairs and 27 test structures.

Besides adding H$_2$ and applying the shared adsorption rules, the catalyst version changes only the prompt. It marks the slab and slab+H pairs, states that the per element constants cancel in $\Delta E_\mathrm{H}$ and gives the DFT and BAM values of $\Delta E_\mathrm{H}$ for each example pair. The prompt and the range check were designed after test results of earlier versions had been seen, and three of the five Claude runs come from the three repeat study described in Section~S2.2.

\subsection{Metal organic frameworks}
\label{sec:mof-methods}

The MOF evaluation uses 33 frameworks from the filtered ODAC25 validation split.\cite{sriram2025odac25} Only pristine frameworks, the original experimental structures that ODAC25 takes from ODAC23, are used, without the variants with added defects, functional groups or generated structures, and further filters are described in Section~S2.3. The reference energies are PBE with D3(BJ) dispersion and the ODAC25 $k$ point correction. Adsorption energies follow Eq.~\eqref{eq:adsorption} with $\nu=1$. The isolated CO$_2$ molecule is not in ODAC25, so its DFT energy is inferred from the stored adsorption energies.

Each framework has several relaxation trajectories that start from different CO$_2$ placements. About one third are set aside, and the final relaxed frame of one of them is the test structure of that framework. The examples are 20 frames spaced along the other trajectories, including intermediate frames, the lowest energy bare framework outside the set aside trajectories and the isolated CO$_2$ molecule. Each framework gets its own correction.

In addition to the shared adsorption rules, the MOF version differs in two respects. First, the prompt describes the examples as structures of a single framework whose DFT reference includes D3 dispersion. Second, because BAM was trained without a dispersion correction, the term library includes a damped pairwise dispersion term. For a chosen element pair, this term contributes $-C_6/(r^6+R_0^6)$ for each atom pair at distance $r$, which decays as $r^{-6}$ at large $r$ and remains finite at short $r$. Like the other pair terms, it is switched off at 5~\AA\ and included in the range check. Its parameter ranges are listed in Table~S1.

The main comparison uses 20 example frames with the dispersion term. This term was introduced after the version without it had been evaluated on the test structures. The comparison of example set sizes uses the same version with 5, 10 and 20 frames (Section~S2.3).

\subsection{CF$_4$ on alumina surfaces}
\label{sec:cf4-methods}

The CF$_4$ test uses $\gamma$-Al$_2$O$_3$(110), $\gamma$-Al$_2$O$_3$(100) and $\alpha$-Al$_2$O$_3$(001) slabs built from the structures of Luo et al.\cite{luo2024cf4} Adsorption energies follow Eq.~\eqref{eq:adsorption}, and unlike the original study, the DFT reference uses PBE without dispersion.

A separate correction is written for each surface from its own examples and applied to all test sites of that surface. The examples of each surface are 10, 8 or 6 frames with a CF$_4$ molecule placed away from the test sites, together with the bare slab and the molecule (Fig.~\ref{fig:cf4-structures} and Table~S\ifbulk 13\else 11\fi). The shared adsorption rules and the four common terms are used unchanged. A comparison of example set sizes uses 3 and 5 frames per surface, drawn at random from these, for both the written correction and fine tuning (Section~S2.4).

\ifbulk
\section{Results and Discussion: Battery and Bulk Materials}

\subsection{Energy and force errors}

\begin{table}[t]
\caption{Energy and force errors on 50 queries per test set. Intervals are paired 95\% bootstrap intervals of the improvement over BAM (BAM error minus corrected error). The last column counts queries whose energy and force errors are both below those of BAM.}
\label{tab:main}
\footnotesize
\setlength{\tabcolsep}{3.5pt}
\resizebox{\textwidth}{!}{%
\begin{tabular}{lrrrrrrrrr}
\toprule
 & & \multicolumn{3}{c}{Energy MAE (meV/atom)} & \multicolumn{4}{c}{Force MAE (eV/\AA)} & \\
\cmidrule(lr){3-5}\cmidrule(lr){6-9}
Test set & Contexts & BAM & Corrected & Improvement & BAM & Corrected & Change & Improvement & Both \\
\midrule
Al alloys & 25 & 42.24 & 13.44 & 21.5--36.8 & 0.0983 & 0.0882 & $-10.2\%$ & 0.0052--0.0151 & 31 \\
LPSCl interpolation & 1 & 59.29 & 2.24 & 56.0--58.0 & 0.1743 & 0.1675 & $-3.9\%$ & 0.0060--0.0075 & 50 \\
LPSCl extrapolation & 1 & 54.82 & 6.91 & 45.6--50.2 & 0.1942 & 0.1835 & $-5.5\%$ & 0.0089--0.0124 & 48 \\
BSW temperature & 7 & 285.09 & 2.96 & 245.9--319.0 & 0.2383 & 0.1947 & $-18.3\%$ & 0.0314--0.0570 & 50 \\
BSW composition & 1 & 276.83 & 48.40 & 228.3--228.6 & 0.2350 & 0.2206 & $-6.1\%$ & 0.0135--0.0153 & 50 \\
BN & 1 & 386.02 & 46.17 & 327.4--350.7 & 1.7875 & 1.6717 & $-6.5\%$ & 0.0801--0.1521 & 40 \\
NMO composition & 16 & 94.69 & 14.75 & 65.6--91.7 & 0.1360 & 0.1207 & $-11.3\%$ & 0.0113--0.0195 & 40 \\
NMO composition and phase & 19 & 94.69 & 15.90 & 65.5--90.1 & 0.1360 & 0.1232 & $-9.4\%$ & 0.0093--0.0163 & 42 \\
\bottomrule
\end{tabular}}
\end{table}

Table~\ref{tab:main} lists the errors. In every test set the LM written correction lowers both the energy and the force error of BAM. The energy error falls by 68--99\%, from 42.24 to 13.44 meV/atom for Al and from 285.09 to 2.96 meV/atom for BSW temperature transfer, and the lowest lower bound of the energy intervals is 21.5 meV/atom. The force error falls by 3.9--18.3\%; the largest reductions are for BSW temperature transfer (0.2383 to 0.1947 eV/\AA) and the two NMO sets, and the smallest for LPSCl interpolation (0.1743 to 0.1675 eV/\AA). Every force interval excludes zero, with the lowest lower bound at 0.0052 eV/\AA\ for Al. No context fell back to BAM.

The gain is not uniform across queries. Both errors are below those of BAM for all 50 queries in LPSCl interpolation and in both BSW sets, but only for 31 Al queries and for 40 and 42 NMO queries. For Al, 34 queries have a lower force error than BAM. BAM has large energy offsets for BN and BSW, which is why the absolute energy reductions there are large; for BN the force error of BAM is also an order of magnitude larger than for the other sets, and the correction lowers it from 1.7875 to 1.6717 eV/\AA.

\subsection{What the written terms are}

\begin{table}[t]
\caption{Selected terms, same formula relative energy error and the effect of refinement. Term counts sum over the selected corrections of all contexts; terms per correction are median (range). The force errors compare the best candidate after the five initial calls with the selected candidate after refinement.}
\label{tab:terms}
\footnotesize
\setlength{\tabcolsep}{4pt}
\resizebox{\textwidth}{!}{%
\begin{tabular}{lrrcrrrr}
\toprule
 & \multicolumn{2}{c}{Selected terms} & & \multicolumn{2}{c}{Relative energy (meV/atom)} & \multicolumn{2}{c}{Force MAE (eV/\AA)} \\
\cmidrule(lr){2-3}\cmidrule(lr){5-6}\cmidrule(lr){7-8}
Test set & Born--Mayer & Morse & Terms per correction & BAM & Corrected & Initial calls & Refined \\
\midrule
Al alloys & 27 & 16 & 1 (1--5) & 11.54 & 12.83 & 0.0922 & 0.0882 \\
LPSCl interpolation & 3 & 0 & 3 & 3.08 & 3.10 & 0.1728 & 0.1675 \\
LPSCl extrapolation & 4 & 0 & 4 & 5.13 & 4.53 & 0.1855 & 0.1835 \\
BSW temperature & 9 & 5 & 2 (1--3) & 2.16 & 1.92 & 0.2165 & 0.1947 \\
BSW composition & 2 & 0 & 2 & 2.34 & 2.37 & 0.2218 & 0.2206 \\
BN & 3 & 1 & 4 & 61.18 & 61.13 & 1.7124 & 1.6717 \\
NMO composition & 10 & 9 & 1 (1--2) & 14.12 & 14.05 & 0.1268 & 0.1207 \\
NMO composition and phase & 11 & 13 & 1 (1--2) & 14.12 & 16.19 & 0.1285 & 0.1232 \\
\bottomrule
\end{tabular}}
\end{table}

The 71 selected corrections contain 113 terms, 69 Born--Mayer and 44 Morse terms (Table~\ref{tab:terms}). Embedded atom terms appeared in 20 of the 1549 scored candidates and bond angle terms in none; neither was selected. Sixty eight of the 69 Born--Mayer terms have $B>0$ and therefore add repulsion, with $\rho$ between 0.1 and 0.5~\AA\ (median 0.3~\AA), so they act mainly on the shortest contacts. The Morse terms have $r_0$ between 0.82 and 1.19 times the first neighbour distance of their pair (median 1.00), stiffness $a$ with median 2.0~\AA$^{-1}$, and $D>0$ in 41 of 44 cases. The LM written corrections thus add short range repulsion and adjust bonds near the observed first neighbour distance. The added repulsion would be consistent with the softening reported for pretrained potentials,\cite{deng2025softening} but these data do not establish that softening is the cause.

The terms change forces much more than they change relative energies. Same formula relative energy errors change by at most 2.07 meV/atom: they decrease for LPSCl extrapolation (5.13 to 4.53) and BSW temperature transfer (2.16 to 1.92), increase for Al (11.54 to 12.83) and NMO composition and phase (14.12 to 16.19), and change by less than 0.1 meV/atom in the other four sets. Since per element constants cancel in this quantity, most of the absolute energy gain in Table~\ref{tab:main} is a correction of the per element reference, while the force gain comes from the terms.

BN illustrates how a correction is built. The residual force table of its prompt has a mean projection of $-2.11$~eV/\AA\ for B--N pairs between 1.0 and 1.25~\AA. Every valid candidate of the five initial calls used a single B--N Born--Mayer term, and the best of them lowered the example force error from 1.689 (BAM) to 1.579 eV/\AA. In the refinement calls the model added a B--N Morse term ($r_0=1.50$~\AA, $D=0.08$~eV, $a=3.0$~\AA$^{-1}$), then a B--B repulsion, and finally a weak N--N repulsion ($B=0.005$~eV) for negative N--N residuals between 1.25 and 2.0~\AA, reaching 1.483 eV/\AA\ on the examples. A stronger N--N repulsion tried in between raised the example force error to 1.789 eV/\AA, above that of BAM, and was not kept.

Refinement matters for forces. Sixty nine of the 71 selected candidates come from refinement calls, and in every test set the refined correction has a lower query force error than the best candidate of the five initial calls, with paired intervals that exclude zero (Table~\ref{tab:terms}). The energy error does not always follow: for LPSCl extrapolation it rises from 4.88 to 6.91 meV/atom and for BN from 39.48 to 46.17 meV/atom, because the selection rule prefers the lowest example force error once both errors are below those of BAM.

\subsection{Validity of the answers}

Of the 710 calls, 19 returned no usable answer: 3 for Al, 1 for BSW temperature transfer and 11 and 4 for the two NMO sets. In addition, 269 candidates within the answers were rejected because they violated the format or a parameter range. Rejected answers and candidates are counted and not resent. The model used a median of 1085 reasoning tokens per call.

\subsection{Composition extrapolation}

BSW composition transfer is the one test set in which every query lies outside the composition span of its examples. All 20 examples have one 153-atom composition, so the examples determine only one combination of the five constants, whereas the queries have 150 atoms. The model, which sees the BAM values, shifted all five constants by the same amount, $-0.2367$~eV relative to BAM, and added two weak Born--Mayer terms for H--Sc and H--Ba pairs. The energy error falls from 276.83 to 48.40 meV/atom, the largest remaining error of the eight sets, and the force error falls by 6.1\%. The equal shift is a choice of the model, guided by the BAM values in the prompt, and not a split determined by the data, so the remaining energy error in this set should not be read as learned extrapolation.

\subsection{Comparison with fine tuning on the same examples}

\begin{figure}[t]
\centering
\includegraphics[width=\linewidth]{fig_icl_ft_ndata_llms_big.pdf}
\caption{Corrections written by four language models and fine tuning (FT) on the same randomly drawn examples, as a function of the number of examples. Top: energy MAE on a logarithmic axis; bottom: force MAE; 564 queries for Al, 200 for each LPSCl set and for BSW temperature transfer, 58 for BSW composition transfer and 1000 for BN and NMO. The dotted line is BAM, which uses no examples. The markers of the language models are shifted horizontally so that close values remain visible. In the Al force panel the FT errors with 10 and 20 examples, 0.306 and 0.294 eV/\AA, lie above the axis range and are drawn at its upper edge with their values. The language models run with different reasoning settings, so the panels do not rank them. Each point is one run with one draw.}
\label{fig:ft}
\end{figure}

\begin{table}[t]
\caption{Errors with ten shared examples per material; LM written is the correction written by Claude Sonnet 5.5. Energy MAE in meV/atom and force MAE in eV/\AA\ on the query sets of Fig.~\ref{fig:ft}.}
\label{tab:ft}
\footnotesize
\setlength{\tabcolsep}{4pt}
\begin{tabular}{lrrrrrr}
\toprule
 & \multicolumn{2}{c}{BAM} & \multicolumn{2}{c}{LM written} & \multicolumn{2}{c}{Fine tuning} \\
\cmidrule(lr){2-3}\cmidrule(lr){4-5}\cmidrule(lr){6-7}
Test set & Energy & Force & Energy & Force & Energy & Force \\
\midrule
Al alloys & 41.95 & 0.1026 & 45.05 & 0.1013 & 393.25 & 0.3058 \\
LPSCl interpolation & 59.38 & 0.1728 & 3.12 & 0.1611 & 1.53 & 0.1069 \\
LPSCl extrapolation & 54.59 & 0.1910 & 6.12 & 0.1801 & 2.07 & 0.1284 \\
BSW temperature & 282.72 & 0.2346 & 1.90 & 0.2202 & 1.37 & 0.1236 \\
BSW composition & 276.83 & 0.2340 & 1.91 & 0.2175 & 1.11 & 0.1246 \\
BN & 382.33 & 1.7854 & 36.74 & 1.6163 & 47.80 & 1.7925 \\
NMO & 99.23 & 0.1244 & 24.25 & 0.1093 & 17.36 & 0.1143 \\
\bottomrule
\end{tabular}
\end{table}

Figure~\ref{fig:ft} and Table~\ref{tab:ft} compare the two methods on identical examples; unless stated otherwise, the LM written correction quoted below is that of Claude Sonnet 5.5. Which method is better depends on the material. For LPSCl, fine tuning gives lower energy and force errors than the LM written correction of every language model at every size. For BSW, fine tuning gives much lower force errors with 20 and 10 examples, 0.1128 and 0.1236 eV/\AA\ for temperature transfer against 0.2192 and 0.2202 eV/\AA, but the energy errors are close: with 20 examples the LM written correction has the lower energy error in both BSW sets, 1.83 and 2.17 against 4.28 and 4.13 meV/atom, the GPT-5.6 Luna correction has a higher energy error than fine tuning for temperature transfer and cannot be distinguished from it for composition transfer, and with ten examples fine tuning has the lower energy error. With ten examples, fine tuning lowers the LPSCl interpolation force error from 0.1728 to 0.1069 eV/\AA, by 38\%, whereas the LM written correction lowers it to 0.1611 eV/\AA, by 7\%; for BSW temperature transfer the reductions are 47\% and 6\%. For LPSCl, fine tuning keeps this lead with five examples, with force errors of 0.1105 and 0.1317 eV/\AA\ against 0.1650 and 0.1838 eV/\AA.

For Al the order is reversed for all four language models because fine tuning fails. It raises the energy error above that of BAM at every size, to 239.78, 393.25 and 428.82 meV/atom with 20, 10 and 5 examples, and with 20 and 10 examples it also roughly triples the force error, to 0.294 and 0.306 eV/\AA. The LM written corrections stay close to BAM: the energy error of the Claude correction cannot be distinguished from that of BAM with 20 and 10 examples and is higher with five, 54.31 against 41.95 meV/atom, and its force error is lower than that of BAM at every size by only 1.0--2.2\%, whereas the Qwen3.8-27B correction does not change the Al force error at any size. A likely cause of the fine tuning failure is the few frames with Fe, Mn or Zn, which enter the refitted element references although the queries contain only Mg, Al and Si, together with overfitting to at most 16 frames; this explanation was not tested separately.

For BN and NMO the comparison is mixed. The LM written correction gives the lower BN force error with 20 and 10 examples, 1.6650 and 1.6163 eV/\AA\ against 1.6961 and 1.7925 eV/\AA, and with ten examples fine tuning is slightly worse than BAM, 1.7854 eV/\AA. With five examples, however, fine tuning gives the lowest BN force error of all runs, 1.5727 eV/\AA. For NMO the LM written correction is better in energy and force with 20 examples (19.17 meV/atom and 0.1031 eV/\AA, against 33.24 meV/atom and 0.1251 eV/\AA); with ten examples fine tuning is better in energy and the LM written correction in force, and with five examples the reverse holds. The other models show the same split: with 20 examples Claude, Gemini and GPT all give a lower BN force error than fine tuning, all four models give a lower NMO force error with 20 and 10 examples, and with five examples fine tuning has the lower BN energy and force errors and the lower NMO force error for every model. Each size is a single draw, so these changes with the number of examples should not be read as trends.

With five examples fine tuning becomes unreliable for BSW: its energy errors rise to 46.11 and 41.87 meV/atom, whereas the LM written correction reaches 13.48 and 9.06 meV/atom, and the Gemini and GPT corrections also have lower energy errors than fine tuning. The five frames cover only two W labels and therefore constrain few combinations of the five element references.

The Claude correction lowers the energy error of BAM in every set except Al and the force error in every set, at every size. The four models differ widely. Qwen3.8-27B leaves the BSW energy error close to that of BAM at every size, between 210.46 and 281.19 meV/atom, and GPT-5.6 Luna reaches 51.04 and 91.81 meV/atom for LPSCl interpolation and NMO with five examples. Because the models run with different reasoning settings and one run each, these differences are not a ranking of the models. Paired 95\% intervals separate the two methods in all comparisons quoted in this section except where they are stated to be indistinguishable.

Fine tuning took 3 to 21 minutes of wall time per material and size on a node with four H100 GPUs, whereas the LM written correction needs ten language model calls per example set and no gradient step. The cathode version uses 35 calls, and one run took about as long as fine tuning on its 21 examples on one A100 GPU (Section~\ref{sec:cathode}). On these examples, fine tuning is the more accurate correction for LPSCl and for the BSW forces; the BSW energy errors of the best LM written corrections are comparable to those of fine tuning; the LM written corrections avoid the failure of fine tuning for Al but improve on BAM there by at most 2.2\% in force; and for BN and NMO the better method changes with the number of examples and with the quantity.

\subsection{Scope of the evidence}

The experiment shows that a language model can write a physically named correction of a frozen network, numbers included, that lowers both the energy and the force errors on these test sets. As a representative case, we compared the written numbers with numbers fitted to the same examples for the same selected terms on the LPSCl test set (Section S2). Relative to the fitted numbers, the written numbers gave a small improvement in energy, 6.91 against 7.81 meV/atom, and no difference in force, 0.1835 against 0.1836 eV/\AA. The main results come from one language model and the fine tuning comparison from four, with one run per test set, and the prompts include the BAM reference values. Selection uses the same examples that the model sees, and several test sets have a single shared context, so the bootstrap intervals describe the queries for the completed calls and not repeated runs. Corrections written independently for each context do not define one validated global potential for dynamics. The fine tuning comparison uses one random draw per material and size and one fixed set of fine tuning settings, so it does not establish how either method depends on the number of examples. No new DFT calculation, stress benchmark or molecular dynamics test is reported. The external model receives the example data and cannot be pinned to a fixed version.

\fi

\clearpage
\section{Results and Discussion: Sodium Ion Cathode}
\label{sec:cathode}

\begin{figure}[!h]
\centering
\includegraphics[width=0.65\linewidth]{fig_cathode_side_hd.pdf}
\caption{P2 Na$_x$MnO$_2$ at $x=5/9$: side view of the lowest energy DFT structure among the 12 at this composition. Mn--O bonds outline the MnO$_6$ octahedra, and the dashed box marks one cell.}
\label{fig:cathode-structures}
\end{figure}
\FloatBarrier

\begin{figure}[!h]
\centering
\includegraphics[width=\linewidth]{fig_p2_hull_voltage_v2.pdf}
\caption{Convex hull and voltage of P2 Na$_x$MnO$_2$. Formation energies of all P2 structures (points) from (a) DFT, (b) BAM, and BAM adapted to the same 21 examples by (c) the LM written correction and (d) fine tuning. Circles mark the examples, black lines the hull of each method and dotted lines the DFT hull. (e) Voltage between neighbouring DFT hull vertices, with the Na metal energy of each method.}
\label{fig:p2hull}
\end{figure}

Figure~\ref{fig:cathode-structures} shows a side view of the lowest-energy DFT structure among the 12 P2 Na$_x$MnO$_2$ structures at $x=5/9$. All four language models lower the formation energy error of BAM, but their voltage errors vary between runs. With 35 calls, the five run mean formation energy MAEs on the 279 P2 evaluation structures are 31.68~meV per formula unit for Claude, 44.02 for Gemini, 54.22 for GPT and 52.37 for Qwen, compared with 59.66 for BAM and 39.13 for fine tuning. The five run mean voltage MAEs are 0.37~V for Claude, 0.24~V for Gemini, 0.31~V for GPT and 0.34~V for Qwen, compared with 0.33~V for BAM and 0.73~V for fine tuning. A candidate passes the selection rule in four of five runs for Claude and Gemini and in one of five for GPT and Qwen, and the other runs keep BAM (Fig.~\ref{fig:p2runs}, Tables~S\ifbulk 4\else 2\fi\ and S\ifbulk 5\else 3\fi\ and Figs.~S2 to S5).

Splitting the error by Na content shows where these differences arise (Table~S\ifbulk 6\else 4\fi). On the Na poor side, $x\le 0.5$, Gemini is close to fine tuning, and GPT and Qwen stay close to BAM, partly because four of their five runs keep BAM. Claude is more accurate there than fine tuning, with a five run mean error at each $x$ of 23 to 47~meV per formula unit, compared with 33 to 63~meV. For $x\ge 0.56$, fine tuning is closer to DFT than Claude, and at $x=0.67$ and 0.72 so is BAM. The lower mean error of Claude than of fine tuning therefore comes from the Na poor side, where the P2 structure changes most with Na content (Section~S\ifbulk 4\else 3\fi).

Figure~\ref{fig:p2hull} shows one Claude run, in which the LM written correction gives a formation energy MAE of 21.08~meV per formula unit and a voltage MAE of 0.15~V. Its mean voltage is 2.63~V, compared with 2.60~V from DFT and 1.89~V from fine tuning. The voltage of fine tuning is too low because its Na metal energy is $-2.19$~eV/atom instead of $-1.45$~eV/atom from DFT, even though Na metal is a training example (Section~S\ifbulk 4\else 3\fi). The correction reduces the tilt of the BAM error, the part of opposite sign on the Na poor and Na rich sides, from 62.5 to 14.9~meV per formula unit. Fine tuning reduces it only to 36.5~meV. The correction overcorrects the depth, the part symmetric about $x=0.5$, from 27.6 to $-18.7$~meV (Section~S\ifbulk 4\else 3\fi). Because the constants $c_Z$ cancel in $\Delta E_\mathrm{f}$, these changes come from the written terms, which are Morse terms for O--O, O--Mn and Na--Na pairs, a Na--Na shell term at 3.05~\AA, embedded atom terms on O and Mn, and an angle term on Na.

The correction is less accurate than both BAM and fine tuning in ranking the structures of the same composition by energy. It identifies the lowest energy structure at 6 of 13 $x$ values, with a mean Spearman correlation of 0.65. BAM and fine tuning identify it at 12, with 0.88 and 0.79. The correction recovers 2 of the 5 intermediate DFT hull vertices, whereas BAM and fine tuning recover all 5. Its $E_{\mathrm{hull}}$ MAE stays close to that of BAM, 31.54 versus 32.70~meV, and above the 25.36~meV of fine tuning. Its force MAE rises from 0.102 to 0.112~eV/\AA, compared with 0.089~eV/\AA\ for fine tuning. No selected correction in any run lowers the $E_{\mathrm{hull}}$ MAE below 31.5~meV or identifies the lowest energy structure at more $x$ values than BAM.

Passing the selection rule does not guarantee a gain in every quantity. With 35 calls, every selected correction lowers the $\Delta E_\mathrm{f}$ MAE, to between 20.91 and 53.80~meV per formula unit. One Claude correction, however, raises the voltage MAE to 0.97~V, compared with 0.33~V for BAM, and this run lifts the Claude mean voltage MAE above that of BAM. Two Gemini corrections pass the energy check on the examples but raise the energy MAE on the evaluation structures from 103 to 119~meV/atom. Five of the ten selected corrections raise the $E_{\mathrm{hull}}$ MAE above the 32.70~meV of BAM, and the Qwen correction also raises the voltage MAE (Table~S\ifbulk 5\else 3\fi).

The call budget decides whether a correction is selected at all. After ten calls, only two Claude runs and one Gemini run have a passing candidate, so 17 of the 20 runs keep BAM. With 35 calls, 10 of the 20 runs select a correction (Fig.~\ref{fig:p2runs}). In the run of Fig.~\ref{fig:p2hull}, the first passing candidate appears at call 11, and later candidates rank better until call 26. The larger budget still takes less time than fine tuning for two of the models. Averaged over five runs, a run with 35 calls takes 15.7~min for Claude and 17.0~min for Gemini, compared with 17.8~min for fine tuning on one A100 GPU. GPT takes 21.0~min and Qwen 30.1~min. Claude thus gives lower $\Delta E_\mathrm{f}$ and voltage errors than fine tuning in less time, although the two run under different conditions (Section~S\ifbulk 4\else 3\fi).

\begin{figure}[!h]
\centering
\includegraphics[width=\linewidth]{fig_p2_llm_runs_v5.pdf}
\caption{Five runs of the LM written correction for each language model on the 279 P2 evaluation structures, using the candidates of the first ten calls (circles, labelled 10) and of all 35 calls (squares, labelled 35): (a) $\Delta E_\mathrm{f}$, (b) $E_{\mathrm{hull}}$ and (c) voltage MAE. Bars mark the mean of five runs, dotted lines BAM and dashed lines fine tuning on the same examples. Runs without a passing candidate keep the BAM values and are spread slightly along the horizontal axis.}
\label{fig:p2runs}
\end{figure}

\FloatBarrier

\clearpage
\section{Results and Discussion: Dual Atom Catalysts}
\label{sec:catalyst}

\begin{figure}[!h]
\centering
\includegraphics[width=\linewidth]{fig_cat_structures_9motifs.pdf}
\caption{Dual atom catalyst M$_2$-N$_x$-C with M a transition metal (Co, Fe, Cr, Ir or Pt), rendered from the relaxed DFT structures. Top: one example structure, with motif 4N/b and one adsorbed H, and the H$_2$ molecule (inset). Bottom: top views of the active site of the slab for three of the ten example motifs and three of the nine test motifs. Each label gives the number of N atoms and, after the slash, their arrangement. The structures shown are those of Co, and the other four metals use the same example and test motifs.}
\label{fig:cat-structures}
\end{figure}
\FloatBarrier

Figure~\ref{fig:cat-structures} shows a Co$_2$--N$_x$--C example structure and representative active-site motifs from the example and test sets. Across five runs, all four language models reduced the metal-averaged hydrogen adsorption energy MAE relative to BAM and fine tuning. Each language model was run five times with each of 5, 10 and 20 examples per metal, giving 60 runs in total. Fine tuning used the same example sets and was run once for each metal and example set size. Both the language model correction and fine tuning used the DFT H$_2$ molecule as an additional example. The unmodified BAM model was evaluated without these additional examples. The adsorption energy evaluation used 45 test pairs across five metals and nine nitrogen coordination motifs. With 20 examples per metal, the mean of the five metal-averaged $\Delta E_\mathrm{H}$ MAEs was 0.220~eV for Claude, 0.253~eV for Gemini, 0.243~eV for GPT and 0.342~eV for Qwen, compared with 0.428~eV for BAM and 0.406~eV for fine tuning. Figure~\ref{fig:cat-llms} compares the four models, while Figure~\ref{fig:cat-parity} shows the Claude results, and Tables~S\ifbulk 9\else 7\fi\ and S\ifbulk 10\else 8\fi\ give the errors for each metal. Results for all models are given in Table~S\ifbulk 7\else 5\fi\ and Figs.~S7--S9.  \par Unlike the adsorption energy, which is calculated from the energies of a slab, its H-adsorbed counterpart and H$_2$, the total energy error is evaluated separately for each of the 135 test structures as the absolute difference from DFT divided by the number of atoms. For Claude, the mean energy error was lower than that of BAM. The mean force error remained close to BAM, while fine tuning gave lower force errors. Of the 135 test structures, 23 single-point structures were calculated without the Hubbard $U$ used for the examples and relaxed structures. These single points are not used for the adsorption energy evaluation. For Cr with 20 examples, the energy MAE of 89.9~meV/atom includes one run that retained BAM, for which the error was 380.4~meV/atom. \par \begin{figure}[t] \centering \includegraphics[width=\textwidth]{fig_cat_4llm_horizontal.pdf} \caption{Hydrogen adsorption energy MAE for four language models using (a) 5, (b) 10 and (c) 20 examples per metal. Each point is one run, and the horizontal bar is the mean of five runs. Each run is the equal-weight average of the five metal-specific MAEs, evaluated on 45 test pairs in total (9 for each metal). Dotted and dashed lines show BAM and fine tuning, respectively; fine tuning uses one run per metal and example-set size. Metal-specific results that retained BAM are included in the averages.} \label{fig:cat-llms} \end{figure}

The language models use common instructions with metal-specific examples to propose corrections to BAM. The selected corrections contain different terms for different metals, while runs without a passing candidate retain BAM. Most selected terms are pair terms. H--H terms appear for all five metals and are particularly frequent for Co, Fe and Pt, while metal--H terms also appear for all five metals. Both can appear in the same correction. The selected terms for each model and metal are listed in Table~S\ifbulk 8\else 6\fi. 





The result varies between runs and with the number of examples. More examples do not lower the error steadily. Qwen, for example, gives 0.349~eV with 5 examples, 0.367~eV with 10 and 0.342~eV with 20. Not every metal improves in every run. With 5 examples, the five run mean for Cr is 0.311~eV, above the 0.297~eV of BAM. For Ir with 20 examples, the five Claude runs give 0.215, 0.391, 0.219, 0.366 and 0.388~eV.

\begin{figure}[t]
\centering
\includegraphics[width=\linewidth]{fig_cat_parity_45pairs_Claude_revised.pdf}
\caption{Hydrogen adsorption energies of 45 test pairs against DFT. Panel (a) shows BAM. Panels (b) to (d) show fine tuning with 5, 10 and 20 examples per metal, respectively. Panels (e) to (g) show all five Claude runs with 5, 10 and 20 examples per metal, respectively. Each bottom-row panel contains five predictions for each of the same 45 pairs, not 225 independent pairs. MAE is averaged over the individual pair predictions without averaging predictions across runs. Each pair has equal weight, which is equivalent to equal metal weighting because each metal has nine pairs. The Cr run that retained BAM is included.}
\label{fig:cat-parity}
\end{figure}




\FloatBarrier


\clearpage

\section{Results and Discussion: CO$_2$ Adsorption in MOFs}
\label{sec:mof}

\begin{figure}[!h]
\centering
\includegraphics[width=\linewidth]{fig_mof_structures.pdf}
\caption{CO$_2$ adsorption in the Zn-based MOF FUBNOH, rendered from DFT structures. Top: an example structure showing CO$_2$ and nearby linkers with their bound Zn atoms; the inset shows isolated CO$_2$. Bottom: four of the five example sites and the test site. CO$_2$ is outlined in white.}
\label{fig:mof-structures}
\end{figure}
\FloatBarrier



\begin{figure}[t]
\centering
\includegraphics[width=0.75\linewidth]{fig_mof_llm_runs.pdf}
\caption{Five runs per language model on 33 MOF test structures, using 20 adsorbed structures and ten calls per MOF. (a) Adsorption energy MAE. (b) Number of MOFs with a selected correction; the others retain BAM. Squares show individual runs and bars show their mean. In (a), dotted and dashed lines indicate BAM and fine tuning, respectively. The grey line in (b) marks 33 MOFs.}
\label{fig:mof-llm-runs}
\end{figure}



Figure~\ref{fig:mof-structures} shows example and test structures for CO$_2$ adsorption in one MOF, FUBNOH. With ten calls and 20 examples per MOF, all four language models reduce the adsorption energy MAE relative to BAM and fine tuning in each of five runs (Fig.~\ref{fig:mof-llm-runs}). On the 33 test structures, the mean $E_\mathrm{ads}$ MAE over five runs is 0.162~eV for Claude, 0.165~eV for Gemini, 0.215~eV for GPT and 0.202~eV for Qwen, compared with 0.260~eV for BAM and 0.337~eV for fine tuning.

\begin{figure}[t]
\centering
\includegraphics[width=\linewidth]{fig_mof_parity.pdf}
\caption{CO$_2$ adsorption energies of the 33 test structures against DFT, from BAM (left), fine tuning (middle) and one Claude run of the LM written correction (right). The last two use 20 examples. BAM uses no examples. The MAE in each panel is taken over the 33 structures.}
\label{fig:mof-parity}
\end{figure}



Figure~\ref{fig:mof-parity} shows one Claude run, in which the LM written correction gives an $E_\mathrm{ads}$ MAE of 0.153~eV. Fine tuning has a higher adsorption energy error than BAM, although its mean error for the loaded frameworks at the test sites is only 0.054~eV. Its errors are 0.401~eV for the bare frameworks and 0.300~eV for the isolated CO$_2$ molecule, even though both are training structures. Both energies enter $E_\mathrm{ads}$ directly. Corresponding results for the first Gemini, GPT and Qwen runs are shown in Figs.~S11--S13.

In the Claude run shown in Fig.~\ref{fig:mof-parity}, the selected corrections included a dispersion term for 20 of the 33 MOFs, mainly for H--O and O--O pairs. For ten MOFs, the selected corrections contained Born--Mayer or Morse terms without a dispersion term. For the remaining three MOFs, no candidate passed selection and BAM was retained. An O--O term can also act within CO$_2$, so it cannot be attributed only to dispersion between the adsorbate and the framework. In a separate Claude run without the dispersion term, the adsorption energy MAE was 0.179~eV on the same 33 test structures (Section~S2.3). \par Across the 20 runs, five for each of the four language models, between 2 and 13 MOFs per run had no passing candidate and retained BAM. XIHRIQ retained BAM in 15 of the 20 runs, while LAZJAZ\_SL and XUKLAS each retained it in 11 runs.



With 20 examples, the five runs of each language model differ by at most 0.028~eV, less than the smallest reduction of any run relative to BAM, 0.031~eV. For one Claude run per example set size, the adsorption energy MAE is 0.172, 0.147 and 0.153~eV with 5, 10 and 20 examples, respectively, with no steady change across these sizes (Table~S\ifbulk 12\else 10\fi). Processing times are reported in Fig.~S10.

\FloatBarrier


Each correction is written from 5, 10 or 20 frames of one MOF and evaluated at the final structure of a set aside trajectory of that same framework. The results therefore show adaptation within a framework, not transfer of a common correction across MOFs. Applied unchanged to the test structures of the other 32 frameworks, the corrections of the run of Fig.~\ref{fig:mof-parity} raise the $E_\mathrm{ads}$ MAE to 0.298~eV, compared with 0.260~eV for BAM. The evaluation comprises 33 test structures, one per MOF, with an assumed CO$_2$ geometry and pristine frameworks of at most 200 atoms. Lower adsorption energy errors alone do not establish better forces or a more accurate potential energy surface.

\FloatBarrier

\clearpage
\section{Results and Discussion: CF$_4$ Adsorption on Al$_2$O$_3$}
\label{sec:cf4}

\begin{figure}[!h]
\centering
\includegraphics[width=\linewidth]{fig_cf4_structures.pdf}
\caption{CF$_4$ adsorption on Al$_2$O$_3$ surfaces, rendered from the DFT structures. Top: the three surfaces with CF$_4$ at their first test site after relaxation without the dispersion correction, seen from the side as a section of the periodic slab, and the isolated CF$_4$ molecule (right). Bottom: $\gamma$(110) at four of its ten example frames, in which the rigid CF$_4$ molecule is placed on the bare slab and in the rightmost one an F atom is directed at a surface Al atom, and at its three test sites. Each panel shows a 12~\AA\ wide section of the slab around the CF$_4$, seen from slightly above.}
\label{fig:cf4-structures}
\end{figure}
\FloatBarrier

\subsection{Adsorption energies}

\begin{figure}[t]
\centering
\includegraphics[width=\linewidth]{fig_cf4_eads.pdf}
\caption{CF$_4$ adsorption energies against the DFT reference without dispersion. (a) Error of $E_\mathrm{ads}$, predicted minus DFT, at the seven test sites for BAM, fine tuning and the LM written correction of one Claude run. Zero is the DFT value, and a negative error means that CF$_4$ binds more strongly than in DFT. (b) $E_\mathrm{ads}$ MAE over the seven sites for five runs of each of Claude, Gemini, GPT and Qwen, each with ten calls. Bars mark the mean of five runs, the dotted line BAM and the dashed line fine tuning.}
\label{fig:cf4-eads}
\end{figure}

\begin{table}[t]
\caption{$E_\mathrm{ads}$ MAE (eV) on the seven test structures for each surface, with the number of test sites in parentheses. The example sets have 10 frames for $\gamma$(110), 8 for $\gamma$(100) and 6 for $\alpha$(001). The reference has no dispersion correction. LM written: one Claude run. Fine tuning was run once.}
\label{tab:cf4-ads}
\footnotesize
\setlength{\tabcolsep}{5pt}
\begin{tabular}{lrrrr}
\toprule
Method & $\gamma$(110) (3) & $\gamma$(100) (3) & $\alpha$(001) (1) & All (7) \\
\midrule
BAM & 0.288 & 0.205 & 0.120 & 0.228 \\
Fine tuning & 0.208 & 0.182 & 0.111 & 0.183 \\
LM written & 0.090 & 0.069 & 0.098 & 0.082 \\
\bottomrule
\end{tabular}
\end{table}

\begin{table}[t]
\caption{LM written correction of the CF$_4$ test with five runs per language model, each with ten calls and the example frames of Table~\ref{tab:cf4-ads}. Selected is the mean number of surfaces, of 3, with a passing candidate, and the other surfaces keep BAM. $E_\mathrm{ads}$ MAE in eV over the seven test structures is the mean of five runs, with the range in parentheses. }
\label{tab:cf4-llm}
\footnotesize
\setlength{\tabcolsep}{5pt}
\begin{tabular}{lrl}
\toprule
Model & Selected & $E_\mathrm{ads}$ MAE \\
\midrule
Claude & 3.0 & 0.086 (0.078--0.096) \\
Gemini & 2.8 & 0.104 (0.071--0.168) \\
GPT & 3.0 & 0.106 (0.081--0.152) \\
Qwen & 1.4 & 0.182 (0.147--0.224) \\
\midrule
BAM & & 0.228 \\
Fine tuning & & 0.183 \\
\bottomrule
\end{tabular}
\end{table}

All four language models lower the adsorption energy error of BAM, and the means of Claude, Gemini and GPT also lie below that of fine tuning. Over five runs, the mean $E_\mathrm{ads}$ MAE on the seven test structures is 0.086~eV for Claude, 0.104~eV for Gemini, 0.106~eV for GPT and 0.182~eV for Qwen. BAM gives 0.228~eV and fine tuning 0.183~eV, so the Qwen mean is close to that of fine tuning (Fig.~\ref{fig:cf4-eads}b and Table~\ref{tab:cf4-llm}). Claude and GPT select a correction for all three surfaces in every run, and Gemini for 14 of its 15 surfaces. Over the five runs, Qwen selects a correction for 7 of 15 surfaces and never for $\gamma$(100).

The gain comes from a repulsion between F and Al that weakens the binding predicted by BAM. BAM binds CF$_4$ too strongly at every test site, by 0.12 to 0.58~eV. Of the 45 surfaces in the runs of Claude, Gemini and GPT, 42 select a correction containing an F--Al term. This term also appears in 5 of the 7 corrections that Qwen selects.

Figure~\ref{fig:cf4-eads}a and Table~\ref{tab:cf4-ads} show one Claude run, in which the LM written correction gives an $E_\mathrm{ads}$ MAE of 0.082~eV. It is closer to DFT than both BAM and fine tuning at every test site. Paired 95\% intervals over the seven structures exclude zero for the MAE difference, corrected minus baseline, from BAM ($-0.146$~eV, $[-0.235, -0.082]$) and from fine tuning ($-0.101$~eV, $[-0.150, -0.057]$). The selected corrections of this run contain an F--Al Born--Mayer repulsion on all three surfaces, with an O--F repulsion added on $\gamma$(100). On the example frames, they lower the $E_\mathrm{ads}$ MAE from 0.223 to 0.034~eV on $\gamma$(110) and from 0.292 to 0.066~eV on $\gamma$(100). On $\alpha$(001), the MAE on the example frames falls from 0.295 to 0.113~eV. Fine tuning leaves errors of 0.163~eV on $\gamma$(110), 0.247~eV on $\gamma$(100) and 0.142~eV on $\alpha$(001) on its own example frames. The largest remaining error of the correction is at the threefold coordinated Al site of $\gamma$(110), the only site that binds CF$_4$ in the reference. There the correction gives $-0.370$~eV, compared with $-0.179$~eV from DFT and $-0.754$~eV from BAM. On $\alpha$(001), the test site error changes little, from 0.120 to 0.098~eV.

No correction changes the forces, as the force MAE on the test structures stays between 0.344 and 0.347~eV/\AA\ in all 20 runs, compared with 0.347~eV/\AA\ for BAM. Fine tuning gives the lowest force MAE, 0.268~eV/\AA, but this gain does not carry over to the adsorption energies.

The result also depends on the run and on the language model. The mean errors of the five Claude runs differ by at most 0.017~eV. Paired intervals exclude zero relative to BAM in all 15 runs of Claude, Gemini and GPT, and relative to fine tuning in 13 of them. The two exceptions are a Gemini run with an MAE of 0.168~eV and a GPT run with an MAE of 0.152~eV. In the Gemini run, no candidate passes on $\gamma$(110), which keeps BAM. In the GPT run, the $\gamma$(100) correction contains only a C--Al term. The Qwen runs range from 0.147 to 0.224~eV, and two of them stay within 0.004~eV of BAM (Table~S\ifbulk 14\else 12\fi).

Drawing 5 or 3 random frames per surface instead of all frames changes the Claude result little. Five Claude runs give a mean $E_\mathrm{ads}$ MAE of 0.093~eV (0.085 to 0.105) with 5 frames and 0.095~eV (0.077 to 0.113) with 3 frames. The mean with all frames is 0.086~eV (Table~S\ifbulk 15\else 13\fi). Fine tuning on the same frames gives 0.197~eV with 5 frames and 0.198~eV with 3 frames. In all ten runs, every surface selects a correction with an F--Al term, and paired 95\% intervals exclude zero relative to both BAM and fine tuning. Of these smaller sets, only the 5 frames of $\alpha$(001) contain a frame with an F atom directed at a surface Al atom. These runs were performed after the results with all frames had been seen.

Processing times are reported in Fig.~S14.

\FloatBarrier
\subsection{Scope of the CF$_4$ evidence}

Each correction is written from 3 to 10 frames of one surface and evaluated at other sites on the same surface. The evaluation set of seven test structures is small, with three sites on two surfaces and one on the third. The reference omits the dispersion correction used by Luo et al., so six of the seven sites are nearly nonbinding. The results therefore do not test predictions against their PBE and D3 reference. Six test sites were placed by us and need not coincide with those of Luo et al. The examples are rigid molecular placements on a fixed slab rather than relaxation frames. This version was evaluated after the method applied unchanged had failed on CF$_4$, and some prompt variants were developed after early test feedback. The results are therefore a development stage evaluation, not an independent validation. Per element constants cancel in the adsorption energies, and lower adsorption energy errors alone do not establish better forces or a more accurate potential energy surface.

\FloatBarrier

\section{Conclusions}

\ifbulk
A frozen language model can write a correction of a frozen interatomic network that contains physically named terms, with no fitter between its answer and the prediction. Given 20 labeled examples, a table of residual forces and parameter ranges taken from the examples, it wrote per element constants together with Born--Mayer and Morse terms, and these corrections lowered the energy error by 68--99\% and the force error by 3.9--18.3\% in all eight test sets, with force intervals that exclude zero. The energy gain is mostly a correction of the per element reference, since same formula energy differences change by at most 2.07 meV/atom, whereas the force gain comes from short range repulsion and bond terms near the observed first neighbour distances. On identical randomly drawn examples, full fine tuning of BAM gave lower energy and force errors for LPSCl and lower force errors for BSW, the LM written corrections of four language models avoided the degradation that fine tuning caused for Al while staying at the BAM level, and for BN and NMO neither method was better throughout. The evidence supports force and energy correction by a language model with frozen neural weights on these test sets.
\else
A frozen language model can write a correction of a frozen interatomic network that contains physically named terms, with no fitter between its answer and the prediction.
\fi

For P2 Na$_x$MnO$_2$, a version with 35 calls and a prompt that describes the hull error by its depth and tilt lowered the formation energy error in four of five Claude runs, from 59.66 to 20.91--31.36~meV per formula unit, and the voltage error to 0.15--0.23~V in three of them, while the ordering of structures at the same composition became worse and fine tuning gave the lowest $E_{\mathrm{hull}}$ error. Gemini reached similar formation energy errors in some runs, GPT and Qwen selected a correction in one run of five each, and with the ten calls of the method 17 of 20 runs kept BAM. These results come from one example set.

For dual atom catalysts, the method applied unchanged gave hydrogen adsorption energy errors close to BAM, since the per element constants cancel in this quantity. With the adsorption energy stated as the target and a check on short range forces, the five run Claude mean lowered the adsorption energy MAE on unseen N coordination motifs from 0.428~eV to 0.220--0.247~eV with 5 to 20 examples. The five run mean force MAE stayed within 0.05~eV/\AA\ of BAM across metals and example set sizes. Fine tuning on the same examples gave lower force errors but adsorption energy errors of 0.406--0.509~eV. With 20 examples per metal, the five run means are 0.220, 0.253, 0.243 and 0.342~eV for Claude, Gemini, GPT and Qwen, respectively, all below BAM and fine tuning.

For CO$_2$ adsorption in 33 MOFs, the MOF version with an added dispersion term lowered the adsorption energy error at the test structure of each framework from 0.260 to 0.153~eV with 20 examples, while fine tuning on the same examples raised it to 0.337~eV because the bare framework and the CO$_2$ molecule were not reproduced. Five runs each with Claude, Gemini, GPT and Qwen, all with the ten calls of the method, stayed below both, between 0.147 and 0.229~eV. With 5 and 10 examples instead of 20, one Claude run gave 0.172 and 0.147~eV. A correction written for one framework did not transfer to other frameworks.
For CF$_4$ on three Al$_2$O$_3$ surfaces, the same version without the dispersion term lowered the adsorption energy error at seven unseen sites from 0.228 to 0.082~eV with 6 to 10 example frames per surface, using a repulsion between F and Al on every surface, while fine tuning on the same examples reached 0.183~eV, and five runs each with Claude, Gemini and GPT gave 0.071 to 0.168~eV, whereas Qwen selected a correction for fewer than half of the surfaces and gave 0.147 to 0.224~eV. With 3 or 5 random frames per surface, ten Claude runs gave 0.077 to 0.113~eV. These results come from seven test structures and a reference without dispersion, and this version was run after the method applied unchanged had failed on this system.

\begin{suppinfo}
Section S1: the common prompt, term parameter ranges (Table S1), model execution settings and response handling. Section S2: data preparation, additional correction terms, prompt changes, selection criteria and comparison settings for each application.
\ifbulk
Section S3: comparison of the numbers written by the language model with numbers fitted to the same examples for the same selected terms on the LPSCl test set (Tables S2 and S3). Section S4: all runs of the cathode version and the error at each Na content (Tables S4 to S6 and Figs.~S2 to S5). Section S5: the five catalyst runs, selected terms and metal-specific errors (Tables S7 to S10). Section S6: the 33 MOFs with their metals, trajectories and test structures and the example set size comparison (Tables S11 and S12). Section S7: the CF$_4$ example frames and test structures, additional reference checks and every run of the CF$_4$ version (Tables S13 to S15).
\else
Section S3: all runs of the cathode version and the error at each Na content (Tables S2 to S4 and Figs.~S2 to S5). Section S4: the five catalyst runs, selected terms and metal-specific errors (Tables S5 to S8). Section S5: the 33 MOFs with their metals, trajectories and test structures and the example set size comparison (Tables S9 and S10). Section S6: the CF$_4$ example frames and test structures, additional reference checks and every run of the CF$_4$ version (Tables S11 to S13).
\fi
\end{suppinfo}

\bibliography{references}

\end{document}
